\documentclass[11pt]{article}

\usepackage[a4paper,margin=1in]{geometry}
\usepackage{amsmath,amssymb,amsthm}
\usepackage{booktabs}
\usepackage{array}
\usepackage{longtable}
\usepackage{xcolor}
\usepackage{listings}
\usepackage{graphicx}
\usepackage[colorlinks=true,linkcolor=blue!55!black,citecolor=blue!55!black,
            urlcolor=blue!55!black]{hyperref}

\allowdisplaybreaks
\numberwithin{equation}{section}

\newcommand{\E}{\mathbb{E}}
\newcommand{\R}{\mathbb{R}}
\newcommand{\N}{\mathbb{N}}
\newcommand{\one}{\mathbf{1}}
\newcommand{\code}[1]{\texttt{\small #1}}
\newcommand{\D}{\mathrm{D}}
\newcommand{\F}{\mathrm{F}}
\newcommand{\bx}{\mathbf{x}}
\newcommand{\bS}{\mathbf{S}}

\theoremstyle{definition}
\newtheorem{definition}{Definition}[section]
\theoremstyle{plain}
\newtheorem{proposition}{Proposition}[section]
\theoremstyle{remark}
\newtheorem{remark}{Remark}[section]

% ---- code listing style -------------------------------------------------
\definecolor{kwc}{RGB}{0,90,150}
\definecolor{cmc}{RGB}{110,110,110}
\definecolor{stc}{RGB}{140,60,20}
\definecolor{frc}{RGB}{200,200,200}
\definecolor{bgc}{RGB}{250,250,248}

\lstdefinestyle{py}{
  language=Python,
  basicstyle=\ttfamily\scriptsize,
  keywordstyle=\color{kwc}\bfseries,
  commentstyle=\color{cmc}\itshape,
  stringstyle=\color{stc},
  numbers=left,
  numberstyle=\tiny\color{cmc},
  numbersep=7pt,
  breaklines=true,
  breakatwhitespace=true,
  showstringspaces=false,
  frame=single,
  rulecolor=\color{frc},
  backgroundcolor=\color{bgc},
  framesep=5pt,
  xleftmargin=16pt,
  captionpos=b,
  columns=fullflexible,
  keepspaces=true,
  aboveskip=8pt,
  belowskip=8pt,
}
\lstset{style=py}
\renewcommand{\lstlistingname}{Listing}

% ---- editorial callouts -------------------------------------------------
\definecolor{notec}{RGB}{120,60,20}
\newcommand{\trick}[1]{\par\smallskip\noindent
  \textcolor{notec}{$\blacktriangleright$}\;\emph{#1}\par\smallskip}
\newcommand{\warn}[1]{\par\smallskip\noindent
  \textcolor{red!60!black}{$\bigstar$}\;\emph{#1}\par\smallskip}

\title{\vspace{-1.4cm}\bfseries A Computational Manual for the\\
Global Projection Solution\\[6pt]
\large Chebyshev--Smolyak projection methods, quadrature over a discrete
default fork,\\ and the solver architecture of \code{code/global/solver\_recursive/}}
\author{Computational appendix}
\date{\today}

\begin{document}
\maketitle

\begin{abstract}
\noindent This manual explains, from first principles and in full detail, how the
two-country monetary-union model with sovereign risk and constrained financial
intermediaries is solved \emph{globally} by a Chebyshev--Smolyak projection method.
It is written for economists who are comfortable with structural macroeconomics ---
Euler equations, Bellman equations, market clearing --- but who have never
implemented a non-linear global solution. Part~1 develops the mathematics of
projection methods: what a projection method \emph{is}, why Chebyshev polynomials
and Chebyshev-extrema nodes rather than any other basis and any other points, and
how Smolyak sparse grids defeat the curse of dimensionality in a seven-dimensional
state space. Part~2 shows how the model's theoretical equations become a concrete
residual function $R(\bx;\bS,\Gamma)=0$, and how expectations over a continuous risk
factor \emph{and} a discrete default fork are computed by Gauss--Hermite quadrature.
Part~3 maps every mathematical object onto the exact file, function and variable
that implements it, with code excerpts. Part~4 --- the longest --- documents the
``under the hood'' machinery: time iteration, damping, homotopy continuation, frozen
continuations, log-space collocation, smoothed guards, failed-point masking, the
Newton solver and its Jacobian, vectorisation, JIT compilation, and how solution
accuracy is measured. Throughout, choices that were tried and \emph{rejected} are
documented alongside those that were kept, because in this class of problem the
failures are as informative as the successes.

\medskip
\noindent\emph{Companion documents.} The file \code{model\_equations.pdf} states
the structural environment: preferences, technology, the intermediary problem, the
government, and the equilibrium definition. The note
\code{global\_solver\_methods.pdf} is a gentler and less complete introduction to
the same numerical ideas; this manual is the technical reference and supersedes it
where the two overlap.
\end{abstract}

\clearpage
\tableofcontents
\clearpage

%%=====================================================================
\section*{Part 0. Orientation}
\addcontentsline{toc}{section}{Part 0. Orientation}
%%=====================================================================
\setcounter{section}{0}

\section{What is being solved, and why globally}

\subsection{The object of interest}

The model is a recursive competitive equilibrium in which every equilibrium quantity
is a \emph{function} of an aggregate state. Let
\begin{equation}
\bS_t=\bigl(K_{\D,t},\,K_{\F,t},\,P_{\D,t},\,P_{\F,t},\,B_{\D,t},\,s_t,\,Z_{\D,t}\bigr)\in\R^{7},
\qquad d_t\in\{0,1\},
\label{eq:state}
\end{equation}
collect the two capital stocks, the two intermediaries' gross deposit obligations,
the risky sovereign debt stock, the latent sovereign-risk factor and total factor
productivity, together with a discrete default regime $d$. We must find functions
\begin{equation}
\bx(d,\bS)=\bigl(N_\D,\;N_\F,\;K'_\D,\;K'_\F,\;r_\D,\;r_\F,\;p\bigr)(d,\bS)
\label{eq:policy-vector}
\end{equation}
--- hours in each country, next period's capital in each country, the two deposit
rates and the terms of trade --- along with valuation functions
$\varphi_X(d,\bS)$ (the intermediary's marginal value of net worth), $q^{X}(d,\bS)$
(sovereign bond prices), and household aggregates $C_X(d,\bS)$, $A_X(d,\bS)$, such
that at \emph{every} point of the state space the seven equilibrium conditions of
Table~\ref{tab:seven} hold.

The unknowns are therefore not numbers but \textbf{functions on $\R^{7}\times\{0,1\}$}.
Since a computer cannot store an arbitrary function, the first task of any global
method is to replace each unknown function by a finite list of coefficients. That is
what a projection method does.

\subsection{Why not linearise?}

The overwhelmingly standard alternative --- perturbation, i.e.\ a Taylor expansion of
the equilibrium conditions around the deterministic steady state --- is
\emph{invalid} for this model, for three independent reasons. Each is worth stating
carefully, because each is the reason a reader from a linear-DSGE background should
expect the machinery in this manual to be necessary rather than ornamental.

\paragraph{(i) The incentive constraint is occasionally binding.} The intermediary's
leverage constraint $\varphi_{X,t}n_{X,t}\ge\lambda_X\mathcal{A}_{X,t}$ holds with
complementary slackness,
\begin{equation}
\mu_{X,t}\ge0,\qquad
\mathrm{slack}_{X,t}\ge0,\qquad
\mu_{X,t}\cdot\mathrm{slack}_{X,t}=0 .
\label{eq:kkt}
\end{equation}
The equilibrium map is therefore only piecewise smooth: it has a \emph{kink} on the
manifold $\{\mu=0\}$. A Taylor expansion around a point on one side of the kink
extrapolates the wrong branch on the other side, and no order of expansion repairs
this, because the map is not analytic there.

\paragraph{(ii) The mechanism of interest is a risk premium, which is identically
zero at first order.} Certainty equivalence holds in a first-order approximation:
the linearised decision rules of a model with $\sigma_s=0.63$ are numerically
identical to those with $\sigma_s=0$. But the entire object this model exists to
measure --- the pass-through of sovereign risk --- \emph{is} a risk premium. In the
bond-pricing condition
\begin{equation}
q^{\D}_t=\frac{\E_t\bigl[\Omega_{t,t+1}\Xi^{\D}_{t+1}\bigr]}
              {\E_t[\Omega_{t,t+1}]R_{\D,t}+\lambda_\D\mu_{\D,t}}
       =\frac{\E_t[\Xi^{\D}_{t+1}]}{R_{\D,t}}
       +\frac{\mathrm{Cov}_t(\Omega_{t,t+1},\Xi^{\D}_{t+1})}{\E_t[\Omega_{t,t+1}]R_{\D,t}}
       -(\text{constraint discount}),
\label{eq:bondprice}
\end{equation}
the covariance term is second order in the shock and vanishes under certainty
equivalence. A first-order solution would price the risky sovereign bond at its
actuarially fair value and report a zero risk premium --- i.e.\ it would return
exactly nothing.

\paragraph{(iii) The relevant experiment is a large, rare, discrete event.} The
default fork is a jump to a state in which the sovereign claim is written down to
$45\%$ of face value. The economy in that branch is far from the steady state:
measured on the solved rules, the intermediary's marginal value of net worth
$\varphi$ reaches more than thirteen times its steady-state level in the deep
default corners. A local approximation is not merely inaccurate there; it is
qualitatively wrong, and --- as \S\ref{sec:mu2} documents --- getting the sign of the
output response right depends on representing that region.

\trick{These three reasons compound. A first-order solution gives no premium; a
second- or third-order perturbation gives a premium but still smooths the
occasionally-binding kink and still expands around a point far from the default
branch. A global method is the only approach that represents both the kink and the
distant branch --- which is why the entire perfect-foresight and
representative-branch machinery this project began with was deleted in favour of the projection
solver.}

\subsection{Road map}

\begin{itemize}
\item \textbf{Part 1} (\S\ref{sec:projection}--\S\ref{sec:box}) --- projection
      methods, Chebyshev bases, Smolyak grids, and the collocation box.
\item \textbf{Part 2} (\S\ref{sec:residual}--\S\ref{sec:quadrature}) --- the residual
      function and numerical integration over the risk factor and the default fork.
\item \textbf{Part 3} (\S\ref{sec:architecture}--\S\ref{sec:mapping}) --- code
      architecture and the mathematics-to-code map.
\item \textbf{Part 4} (\S\ref{sec:timeiter}--\S\ref{sec:accuracy}) --- the solver
      algorithm and every computational trick it relies on.
\end{itemize}

%%=====================================================================
\clearpage
\section*{Part 1. Mathematical foundations of the solution method}
\addcontentsline{toc}{section}{Part 1. Mathematical foundations}
%%=====================================================================

\section{Projection methods from first principles}
\label{sec:projection}

\subsection{The general problem}

Every global solution method starts from the same abstract statement. We seek an
unknown function $g:\mathcal{D}\to\R^{n}$, defined on a domain
$\mathcal{D}\subseteq\R^{d}$, that satisfies a functional equation
\begin{equation}
\mathcal{H}\bigl[g\bigr](\bS)\;=\;0
\qquad\text{for all }\bS\in\mathcal{D},
\label{eq:functional}
\end{equation}
where $\mathcal{H}$ is an operator that may involve $g$ evaluated at other points
(through the transition law) and expectations over future shocks. In our case
$\mathcal{H}$ is the stack of seven equilibrium conditions and $g$ is the collection
of decision rules \eqref{eq:policy-vector}.

Equation~\eqref{eq:functional} is an infinite system: one vector equation for each
of the uncountably many points $\bS\in\mathcal{D}$. A \textbf{projection method}
makes it finite in two steps.

\paragraph{Step 1: parameterise.} Choose a family of basis functions
$\{\Psi_1,\dots,\Psi_n\}$ and approximate
\begin{equation}
g(\bS)\;\approx\;\hat g(\bS;c)\;=\;\sum_{b=1}^{n}c_b\,\Psi_b(\bS),
\label{eq:parameterisation}
\end{equation}
so the unknown function is replaced by the unknown coefficient vector
$c\in\R^{n}$.

\paragraph{Step 2: project.} Substituting \eqref{eq:parameterisation} into
\eqref{eq:functional} defines the \emph{residual function}
\begin{equation}
R(\bS;c)\;\equiv\;\mathcal{H}\bigl[\hat g(\cdot\,;c)\bigr](\bS),
\label{eq:residual-def}
\end{equation}
which will not be identically zero for any finite $n$. We instead require it to be
zero in a weighted-average sense: choose $n$ weighting functions $w_1,\dots,w_n$ and
impose
\begin{equation}
\int_{\mathcal{D}}w_b(\bS)\,R(\bS;c)\,d\bS\;=\;0,
\qquad b=1,\dots,n .
\label{eq:projection-conditions}
\end{equation}
This is $n$ equations in $n$ unknowns.

Different choices of $w_b$ give the classical variants: Galerkin ($w_b=\Psi_b$),
least squares ($w_b=\partial R/\partial c_b$), and --- the one used here ---
\textbf{collocation}, where $w_b(\bS)=\delta(\bS-\bS_b)$ is a Dirac mass at a
selected node $\bS_b$. Collocation reduces \eqref{eq:projection-conditions} to the
simple and intuitive requirement
\begin{equation}
\boxed{\;R(\bS_b;c)=0,\qquad b=1,\dots,n\;}
\label{eq:collocation}
\end{equation}
\emph{the residual is exactly zero at $n$ chosen points of the state space.} Between
those points the residual is not zero; how large it is there is precisely the
accuracy question of \S\ref{sec:accuracy}.

\begin{remark}[Why collocation]
Galerkin projection requires integrating the residual against each basis function
over a seven-dimensional domain, and each residual evaluation here costs a
seven-equation non-linear solve with a multi-branch quadrature inside it. Collocation
requires the residual only at $n$ points. With $n=259$ nodes and two regimes, the
cost difference is the difference between a solve that runs in ninety minutes and one
that does not run at all.
\end{remark}

\subsection{What ``choose a basis'' means in practice}

Two decisions remain: \emph{which} basis functions $\Psi_b$, and \emph{which}
collocation points $\bS_b$. These are not independent --- a bad pairing produces a
near-singular collocation matrix or a wildly oscillating interpolant --- and the next
two sections address them together. The answer used here is: tensor products of
Chebyshev polynomials, evaluated at Chebyshev-extrema nodes, both restricted to a
Smolyak sparse index set.

\section{Chebyshev approximation in one dimension}
\label{sec:chebyshev}

\subsection{Definition and the recurrence}

The Chebyshev polynomials of the first kind are defined on $[-1,1]$ by
\begin{equation}
T_k(x)\;=\;\cos\bigl(k\arccos x\bigr),
\qquad k=0,1,2,\dots
\label{eq:cheb-def}
\end{equation}
The trigonometric definition makes several properties immediate: $|T_k(x)|\le1$ on
$[-1,1]$; $T_k$ has exactly $k$ roots and $k+1$ extrema in $[-1,1]$; and $T_k$
oscillates between $+1$ and $-1$ with equal amplitude --- the equioscillation
property that underlies its optimality. Despite the cosine, $T_k$ \emph{is} a
polynomial of degree $k$, as the three-term recurrence shows:
\begin{equation}
T_0(x)=1,\qquad T_1(x)=x,\qquad
T_{k+1}(x)=2x\,T_k(x)-T_{k-1}(x).
\label{eq:cheb-recurrence}
\end{equation}
Equation~\eqref{eq:cheb-recurrence} --- not \eqref{eq:cheb-def} --- is what the code
evaluates, because it is $O(k)$ per point, numerically stable, and free of
transcendental function calls. It is implemented verbatim:

\begin{lstlisting}[caption={The Chebyshev recurrence. \code{state\_grid.py}, \code{chebyshev\_basis\_1d}.}]
def chebyshev_basis_1d(x, max_deg):
    # T_0..T_max_deg AT POINTS x VIA THE RECURRENCE (SHAPE (len(x), max_deg+1)).
    x = np.asarray(x, dtype=float)
    T = np.empty((x.size, max_deg + 1))
    T[:, 0] = 1.0
    if max_deg >= 1:
        T[:, 1] = x
    for k in range(2, max_deg + 1):
        T[:, k] = 2.0 * x * T[:, k - 1] - T[:, k - 2]
    return T
\end{lstlisting}

The polynomials are orthogonal on $[-1,1]$ with respect to the weight
$(1-x^{2})^{-1/2}$:
\begin{equation}
\int_{-1}^{1}\frac{T_j(x)T_k(x)}{\sqrt{1-x^{2}}}\,dx
=\begin{cases}
0, & j\ne k,\\
\pi, & j=k=0,\\
\pi/2, & j=k\ge1 .
\end{cases}
\label{eq:cheb-orthogonality}
\end{equation}
Orthogonality is what makes the coefficients of a Chebyshev expansion decay rapidly
and independently, so that truncating the series at degree $n$ discards a
well-controlled remainder rather than redistributing error across the retained terms.

\subsection{Why Chebyshev: three properties that matter}

\paragraph{(a) Near-minimax accuracy.} Let $f$ be continuous on $[-1,1]$ and let
$p^{\ast}_n$ be its best degree-$n$ polynomial approximation in the sup norm. The
degree-$n$ Chebyshev interpolant $p_n$ satisfies
\begin{equation}
\|f-p_n\|_{\infty}\;\le\;\Bigl(1+\Lambda_n\Bigr)\,\|f-p^{\ast}_n\|_{\infty},
\qquad
\Lambda_n\;\sim\;\frac{2}{\pi}\log n ,
\label{eq:lebesgue}
\end{equation}
where $\Lambda_n$ is the Lebesgue constant of the node set. The logarithmic growth is
the crucial fact: at $n=64$ the Chebyshev interpolant is within a factor of about
$3.6$ of the unimprovable approximation. For \emph{equispaced} nodes the Lebesgue
constant instead grows like $2^{n}/(e\,n\log n)$ --- exponentially --- which is the
Runge phenomenon of \S\ref{sec:runge}.

\paragraph{(b) Spectral (geometric) convergence for smooth functions.} If $f$ is
analytic in the Bernstein ellipse $E_\rho$ (the ellipse in the complex plane with
foci $\pm1$ and semi-axis sum $\rho>1$), then
\begin{equation}
\|f-p_n\|_{\infty}\;\le\;\frac{2M\rho^{-n}}{\rho-1},
\qquad M=\max_{z\in E_\rho}|f(z)| ,
\label{eq:spectral}
\end{equation}
i.e.\ the error falls \emph{geometrically} in the degree. Adding one degree
multiplies the accuracy by a constant factor, rather than adding a fixed number of
digits as in a finite-difference scheme. If instead $f$ has only $k$ continuous
derivatives, convergence degrades to the algebraic rate $O(n^{-k})$.

\warn{Property (b) is exactly what the occasionally-binding constraint threatens.
The equilibrium map is smooth on each side of $\{\mu=0\}$ but only $C^{0}$ across it.
Where the kink lies inside the collocation box, the geometric rate
\eqref{eq:spectral} is lost and the interpolant develops Gibbs-type oscillations.
Section~\ref{sec:mu2} shows this is not hypothetical: it is the reason the solve
converges at approximation level $\ell=1$ and fails at $\ell=2$.}

\paragraph{(c) A well-conditioned collocation matrix.} Because the basis is
orthogonal and the nodes are the natural ones for that basis, the matrix $\Phi$ of
\S\ref{sec:fitting} is far better conditioned than the Vandermonde matrix of a
monomial basis $\{1,x,x^{2},\dots\}$, whose condition number grows exponentially in
degree. In double precision a monomial basis becomes unusable above degree ten or so;
the Chebyshev basis remains usable into the hundreds.

\subsection{Node placement and the Runge disaster}
\label{sec:runge}

Interpolation error at any $x$ factors as
\begin{equation}
f(x)-p_n(x)\;=\;\frac{f^{(n+1)}(\xi)}{(n+1)!}\prod_{j=0}^{n}\bigl(x-x_j\bigr),
\label{eq:interp-error}
\end{equation}
so the node placement enters only through the \emph{node polynomial}
$\prod_j(x-x_j)$. Minimising its sup norm over node choices is a classical problem
whose solution clusters points towards the endpoints with the arcsine density
$1/(\pi\sqrt{1-x^{2}})$. Equispaced nodes leave the node polynomial exponentially
large near $\pm1$, and interpolants of even innocuous functions --- Runge's
$f(x)=1/(1+25x^{2})$ is the textbook case --- \emph{diverge} as $n$ grows.

This code uses the \textbf{Chebyshev extrema} (Gauss--Lobatto) points,
\begin{equation}
x_j\;=\;-\cos\!\left(\frac{\pi j}{m-1}\right),
\qquad j=0,1,\dots,m-1,
\label{eq:cheb-extrema}
\end{equation}
which include the endpoints $\pm1$ and satisfy the arcsine clustering. The choice of
extrema rather than \emph{roots} (Gauss points) is deliberate and is what makes the
Smolyak construction of \S\ref{sec:smolyak} possible: extrema sets are
\textbf{nested}, so that the level-$i$ point set is a subset of the level-$(i{+}1)$
set. Root sets are not.

\begin{lstlisting}[caption={Nested Chebyshev-extrema levels. \code{state\_grid.py}.}]
def _level_points(i):
    # FULL 1D CHEBYSHEV-EXTREMA SET OF LEVEL i (m = 1, 3, 5, 9, ... POINTS).
    if i == 1:
        return np.array([0.0])
    m = 2 ** (i - 1) + 1
    return -np.cos(np.pi * np.arange(m) / (m - 1))

def _new_points(i):
    # POINTS INTRODUCED AT LEVEL i (DISJOINT ACROSS LEVELS BY NESTEDNESS).
    if i == 1:
        return np.array([0.0])
    if i == 2:
        return np.array([-1.0, 1.0])
    return _level_points(i)[1::2]   # odd positions are absent from level i-1
\end{lstlisting}

The level sequence is $m_i=1,3,5,9,17,\dots$ (i.e.\ $m_1=1$ and
$m_i=2^{\,i-1}+1$ for $i\ge2$), and the ``new points'' at level $i$ are the
odd-indexed members of level $i$, which are exactly those absent from level $i-1$.
Table~\ref{tab:levels} lists the first few levels.

\begin{table}[htbp]
\centering\small
\caption{Nested one-dimensional levels: points and the degrees they introduce.}
\label{tab:levels}
\begin{tabular}{@{}cccl@{}}
\toprule
Level $i$ & $m_i$ (cumulative points) & New points & New degrees \\
\midrule
1 & 1 & $\{0\}$                                  & $\{0\}$ \\
2 & 3 & $\{-1,+1\}$                              & $\{1,2\}$ \\
3 & 5 & $\{-\tfrac{\sqrt2}{2},+\tfrac{\sqrt2}{2}\}$ & $\{3,4\}$ \\
4 & 9 & 4 points                                  & $\{5,6,7,8\}$ \\
5 & 17& 8 points                                  & $\{9,\dots,16\}$ \\
\bottomrule
\end{tabular}
\end{table}

The pairing of ``new points'' with ``new degrees'' in the last column is what keeps
the collocation system square, and is implemented as:

\begin{lstlisting}[caption={Degrees introduced at each level. \code{state\_grid.py}, \code{\_new\_degrees}.}]
def _new_degrees(i):
    # CHEBYSHEV DEGREES INTRODUCED AT LEVEL i (|new degrees| = |new points|).
    if i == 1:
        return np.array([0])
    if i == 2:
        return np.array([1, 2])
    m_prev = 2 ** (i - 2) + 1
    return np.arange(m_prev, 2 ** (i - 1) + 1)
\end{lstlisting}

\subsection{Fitting is a linear solve}
\label{sec:fitting}

Given nodes $\{x_p\}_{p=1}^{n}$ and degrees $\{k_b\}_{b=1}^{n}$, build the
collocation matrix
\begin{equation}
\Phi_{pb}\;=\;T_{k_b}(x_p),
\qquad
\Phi\in\R^{n\times n},
\label{eq:phi-1d}
\end{equation}
and let $y_p=f(x_p)$ be the function values. The interpolation conditions
$\hat f(x_p)=y_p$ are the linear system
\begin{equation}
\Phi\,c\;=\;y
\qquad\Longrightarrow\qquad
c=\Phi^{-1}y .
\label{eq:fit-1d}
\end{equation}
Two implementation points matter. First, $\Phi$ depends only on the \emph{grid}, not
on the function being fitted, so it is built once and LU-factorised once; every
subsequent fit is a triangular back-substitution costing $O(n^{2})$ instead of
$O(n^{3})$. Second, this is what makes it cheap to fit \emph{seventeen} rules in two
regimes on every sweep: they all share one factorisation.

\begin{lstlisting}[caption={One factorisation, many fits. \code{state\_grid.py}, \code{SmolyakGrid}.}]
self._Phi = self._basis_unit(self.points_unit)     # (n, n) collocation basis
self._lu = lu_factor(self._Phi)
...
def fit(self, values):
    # COLLOCATION COEFFICIENTS FROM VALUES AT self.points ((n,) OR (n, k)).
    return lu_solve(self._lu, np.asarray(values, dtype=float))
\end{lstlisting}

\subsection{Never extrapolate}

Equation~\eqref{eq:spectral} holds \emph{inside} the interpolation interval. Outside
it, the Chebyshev interpolant behaves like the leading monomial $x^{k_{\max}}$ and
diverges rapidly: at $x=1.3$ with $k_{\max}=8$, $T_8$ is already of order $10^{3}$.
Every evaluation of the decision rules must therefore be inside the box. The code
enforces this explicitly at the one place where an off-box state can arise --- the
continuation point in the expectation --- by clipping:
\begin{equation}
\bS'_{\text{eval}}=\Pi_{[\text{lo},\text{hi}]}\bigl(\bS'\bigr),
\qquad
\text{code: \code{cont.eval\_all(d\_next, cont.grid.clip(Sn))}} .
\label{eq:clip}
\end{equation}
Clipping is not free --- \S\ref{sec:box} shows that an over-tight box makes the clip
distort the risk process itself --- so the box must be wide enough that clipping is
rare, and how often it binds is reported by the accuracy diagnostic
(\S\ref{sec:accuracy}).

\section{Seven dimensions without the explosion}
\label{sec:smolyak}

\subsection{The curse of dimensionality}

The natural multi-dimensional extension of \S\ref{sec:chebyshev} is the tensor
product: take $m$ nodes per dimension and use all $m^{d}$ combinations, with the
corresponding product basis $\prod_{j}T_{k_j}(x_j)$. In $d=7$ with the modest choice
$m=3$ this is $3^{7}=2{,}187$ points; with $m=5$ it is $78{,}125$; with $m=9$ it is
$4{,}782{,}969$. Each point requires a seven-equation non-linear solve with a
$3M$-node quadrature inside it, in each of two regimes. At roughly a second per
point, $m=5$ alone is a day of computation per sweep, and dozens of sweeps are
needed. Tensor products are unusable here.

\subsection{The Smolyak construction}

Smolyak's (1963) insight, in the nested Chebyshev form of Krueger and Kubler (2004),
is that most of the tensor product is wasted: it spends points on \emph{high-order
interactions} --- terms like $x_1^{8}x_2^{8}x_3^{8}$ --- which contribute almost
nothing to the accuracy of a smooth function. A sparse grid keeps the cheap
combinations and discards the expensive ones, according to a budget on the total
level.

Formally, index each dimension by a level $i_j\ge1$, and define the admissible
multi-index set
\begin{equation}
\mathcal{I}(\ell)\;=\;\Bigl\{\,i=(i_1,\dots,i_d)\in\N^{d}\;:\;
\sum_{j=1}^{d}\bigl(i_j-1\bigr)\;\le\;\ell,
\quad i_j-1\le\ell_j\ \forall j\,\Bigr\},
\label{eq:multi-index}
\end{equation}
where $\ell$ is the \textbf{approximation level} (the code calls it \code{mu}; we
write $\ell$ throughout to avoid collision with the incentive-constraint multiplier
$\mu$) and the optional per-dimension caps $\ell_j$ produce the anisotropic grids of
\S\ref{sec:anisotropic}. The grid and the basis are then unions of small tensor
products over that index set, using the \emph{new} point and degree sets of
Table~\ref{tab:levels}:
\begin{align}
\mathcal{H}(\ell)&=\bigcup_{i\in\mathcal{I}(\ell)}\;
   \Delta_{i_1}\times\Delta_{i_2}\times\cdots\times\Delta_{i_d},
&&\Delta_i=\text{new points at level }i,
\label{eq:smolyak-grid}\\[3pt]
\mathcal{B}(\ell)&=\bigcup_{i\in\mathcal{I}(\ell)}\;
   \mathcal{K}_{i_1}\times\mathcal{K}_{i_2}\times\cdots\times\mathcal{K}_{i_d},
&&\mathcal{K}_i=\text{new degrees at level }i .
\label{eq:smolyak-basis}
\end{align}
Because $|\Delta_i|=|\mathcal{K}_i|$ for every $i$ and the unions are over the same
index set, $|\mathcal{H}(\ell)|=|\mathcal{B}(\ell)|=n$: the collocation matrix is
\textbf{square by construction}, and \eqref{eq:fit-1d} carries over unchanged. This
is the structural reason the ``new points / new degrees'' pairing of
Table~\ref{tab:levels} is set up as it is.

The multi-index enumeration is a depth-first recursion over the remaining level
budget:

\begin{lstlisting}[caption={Enumerating $\mathcal{I}(\ell)$. \code{state\_grid.py}, \code{\_multi\_indices}.}]
def _multi_indices(d, mu, mu_vec):
    # ALL LEVEL MULTI-INDICES i (EACH >= 1) WITH sum(i-1) <= mu AND i-1 <= mu_vec.
    out = []
    def rec(prefix, budget):
        # DEPTH-FIRST ENUMERATION UNDER THE REMAINING LEVEL BUDGET.
        j = len(prefix)
        if j == d:
            out.append(tuple(prefix))
            return
        for lev in range(1, min(budget, mu_vec[j]) + 2):
            rec(prefix + [lev], budget - (lev - 1))
    rec([], mu)
    return out
\end{lstlisting}

and the grid assembly is a loop of small meshgrids:

\begin{lstlisting}[caption={Building the sparse grid and its basis degrees. \code{state\_grid.py}, \code{SmolyakGrid.\_\_init\_\_}.}]
pts, degs = [], []
for i_vec in _multi_indices(self.d, self.mu, self.mu_vec):
    axes_p = [_new_points(i)  for i in i_vec]
    axes_d = [_new_degrees(i) for i in i_vec]
    mesh_p = np.meshgrid(*axes_p, indexing="ij")
    mesh_d = np.meshgrid(*axes_d, indexing="ij")
    pts.append(np.column_stack([m.ravel() for m in mesh_p]))
    degs.append(np.column_stack([m.ravel() for m in mesh_d]))
self.points_unit = np.vstack(pts)
self.degrees     = np.vstack(degs).astype(int)
\end{lstlisting}

The multi-dimensional basis matrix is then the elementwise product of
one-dimensional Chebyshev matrices, one per dimension --- a fully vectorised
operation over all points and all basis functions at once:

\begin{lstlisting}[caption={The tensor basis, vectorised. \code{state\_grid.py}, \code{\_basis\_unit}.}]
def _basis_unit(self, u):
    # BASIS MATRIX AT UNIT-BOX POINTS: PRODUCTS OF PER-DIMENSION CHEBYSHEVS.
    u = np.atleast_2d(u)
    B = np.ones((u.shape[0], self.n))
    for j in range(self.d):
        Tj = chebyshev_basis_1d(u[:, j], self.max_deg)
        B *= Tj[:, self.degrees[:, j]]
    return B
\end{lstlisting}

\subsection{What it buys}

Table~\ref{tab:counts} gives the actual point counts for $d=7$, computed from the
implementation. The savings are dramatic and grow with dimension.

\begin{table}[htbp]
\centering\small
\caption{Smolyak versus tensor-product point counts in $d=7$ (computed from
\code{SmolyakGrid}).}
\label{tab:counts}
\begin{tabular}{@{}lrrrr@{}}
\toprule
Level $\ell$ & Max degree & Smolyak points & Tensor points & Ratio \\
\midrule
1 & 2 & 15  & 2{,}187     & $1:146$ \\
2 & 4 & 113 & 78{,}125    & $1:691$ \\
3 & 8 & 589 & 4{,}782{,}969 & $1:8{,}120$ \\
\midrule
\multicolumn{5}{@{}l@{}}{\emph{Grids actually used:}}\\
$[1,1,1,1,1,3,0]$ (production, risk) & 8 & \textbf{259} & --- & anisotropic \\
$[0,0,1,1,0,3,0]$ (frozen capital)   & 8 & 41  & --- & anisotropic \\
$\ell=1$ isotropic (TFP experiment)  & 2 & 15  & 2{,}187 & $1:146$ \\
\bottomrule
\end{tabular}
\end{table}

The approximation-theoretic guarantee is that the sparse grid loses surprisingly
little. Level $\ell$ reproduces \emph{exactly} every polynomial of total degree
$\le\ell$ in $d$ variables, including all cross terms. At $\ell=2$ this means every
complete quadratic --- all $d$ linear terms, all $d$ squares and all $d(d-1)/2$
interactions --- which is already the entire information content of a second-order
perturbation, but represented globally rather than locally. This property is
regression-tested directly:

\begin{lstlisting}[caption={The exactness property, tested. \code{tests/test\_state\_grid.py}.}]
def test_quadratic_exactness():
    # LEVEL 2 REPRODUCES COMPLETE QUADRATICS INCL. CROSS TERMS (KK04 PROPERTY).
    rng = np.random.default_rng(1); d = 8
    g = SmolyakGrid(-np.ones(d), np.ones(d), mu=2)
    A = rng.standard_normal((d, d)); A = 0.5 * (A + A.T)
    b = rng.standard_normal(d); c0 = rng.standard_normal()
    def f(x):
        return np.einsum("ij,jk,ik->i", x, A, x) + x @ b + c0
    coef = g.fit(f(g.points))
    x = rng.uniform(-1, 1, size=(800, d))
    assert np.max(np.abs(g.eval(coef, x) - f(x))) < 1e-9
\end{lstlisting}

The general error bound (Barthelmann, Novak and Ritter, 2000) for $f$ with bounded
mixed derivatives of order $k$ is
\begin{equation}
\bigl\|f-\hat f_{\ell}\bigr\|_{\infty}
\;=\;O\!\left(n^{-k}\bigl(\log n\bigr)^{(d-1)(k+1)}\right),
\label{eq:smolyak-error}
\end{equation}
against the tensor-product rate $O(n^{-k/d})$ --- the curse of dimensionality is
reduced from an exponent to a logarithmic factor.

\subsection{Anisotropic refinement}
\label{sec:anisotropic}

The isotropic budget in \eqref{eq:multi-index} spends resolution equally on all seven
states. That is wasteful here, because the states are economically very unequal: the
sovereign-risk factor $s$ ranges over $\pm4.35$ in logit units and is where every
non-linearity lives, whereas the two capital stocks move by a few percent over the
ergodic set. Setting per-dimension caps $\ell_j$ --- the argument \code{mu\_vec} ---
concentrates points where the curvature is:
\begin{equation}
\ell_{\text{vec}}=\bigl[\underbrace{1,1}_{K_\D,K_\F},
\underbrace{1,1}_{P_\D,P_\F},\underbrace{1}_{B_\D},
\underbrace{3}_{s},\underbrace{0}_{Z_\D}\bigr]
\;\Longrightarrow\; n=259,\ \ k_{\max}=8 .
\label{eq:mu-vec}
\end{equation}
The risk factor gets degree-8 resolution; TFP, which is never shocked in the risk
experiment, gets a single node (a rule flat in $Z$); the remaining five states get
the minimum non-trivial refinement.

\trick{The anisotropic grid is not a cost-saving convenience --- it changes the
answer. On the isotropic $\ell=1$ grid the long bond reprices by only $-1.3\%$ under
the headline shock, because degree-2 resolution in $s$ cannot represent the convexity
of the price in the risk factor. With $\ell_s=3$ it reprices by $-4.8\%$, the bond
loss crashes intermediary net worth, and the multiplier $\mu$ spikes endogenously.
The mechanism the model exists to measure is invisible at uniform resolution.}

\warn{The opposite error is equally consequential. The variant
$[0,0,1,1,0,3,0]$ (41 points) freezes $K_\D$, $K_\F$ and $B_\D$ at a single node,
making the rules \emph{flat} in capital and debt. This runs in ten minutes instead of
ninety, but it switches off the capital-accumulation channel entirely: measured, the
impact credit-spread response is $+143$\,bp with capital live and $-27$\,bp with it
frozen. A grid choice silently deleted the mechanism and flipped a sign.}

\section{The collocation box}
\label{sec:box}

\subsection{\texorpdfstring{Mapping the economic domain to $[-1,1]^{d}$}
                          {Mapping the economic domain to the unit box}}

Chebyshev theory lives on $[-1,1]^{d}$; the model lives on an economic box
$[\text{lo},\text{hi}]\subset\R^{7}$. The affine map and its inverse are
\begin{equation}
u_j(\bS)=2\,\frac{S_j-\text{lo}_j}{\text{hi}_j-\text{lo}_j}-1,
\qquad
S_j=\text{lo}_j+\tfrac12\bigl(u_j+1\bigr)\bigl(\text{hi}_j-\text{lo}_j\bigr),
\label{eq:box-map}
\end{equation}
implemented as \code{to\_unit} and \code{from\_unit}. The box is centred on the
steady state with per-state bands:
\begin{equation}
\begin{aligned}
K_X&:\ \pm2\%, &\qquad P_\D&:\ \pm12\%, &\qquad P_\F&:\ \pm20\%,\\
B_\D&:\ \pm8\%, &\qquad s&:\ \bar s\pm4.35, &\qquad Z_\D&:\ \pm3\% .
\end{aligned}
\label{eq:bands}
\end{equation}

\subsection{Why these bands, specifically}

Band selection is one of the highest-leverage choices in the whole solve, and the
values in \eqref{eq:bands} are measured rather than guessed.

\paragraph{The $s$ band.} The half-width $4.35$ is Bocola's own coverage, equal to
$\pm2.16$ unconditional standard deviations of the $s$ process
($\sigma_s/\sqrt{1-\rho_s^{2}}=2.02$). This must be checked against the
\emph{conditional} innovation, not just the unconditional spread, because the
continuation states in the quadrature are clipped to the box: if the box is narrow in
conditional standard deviations, the clip mean-reverts $s$ harder than $\rho_s$ does.
An earlier calibration with $\sigma_s=1.5075$ made the box only $1.4$--$2.3$
conditional standard deviations wide, which cut the \emph{effective} persistence of
$s$ from $0.95$ to $0.80$ and removed roughly $62\%$ of the long bond's repricing ---
silently. There is now a regression test for exactly this:

\begin{lstlisting}[caption={Guarding against a box that distorts the shock process. \code{tests/test\_state\_grid.py}.}]
    # and the clip must not distort the process in the ergodic core
    x, w = np.polynomial.hermite_e.hermegauss(5)
    w = w / w.sum()
    grid_s = np.linspace(sp["s_star"] - unc, sp["s_star"] + unc, 21)
    Ec = np.array([w @ np.clip((1 - sp["rho_s"]) * sp["s_star"] + sp["rho_s"] * s
                               + sp["sigma_s"] * x, lo, hi) for s in grid_s])
    rho_eff = np.polyfit(grid_s, Ec, 1)[0]
    assert rho_eff > 0.93, f"box clip cuts effective rho_s to {rho_eff:.3f}"
\end{lstlisting}

\paragraph{The $P_\F$ band.} At the original $\pm8\%$ the sweep degraded from about
the tenth iteration and $22$ of $41$ points stopped clearing permanently, leaving the
fit anchored on frozen cold-start values with a worst residual of $9.6\times10^{-3}$.
At $\pm20\%$ the entire grid clears to $1.1\times10^{-14}$ with zero failures, and
$4.5\times$ faster. The reason is economic: foreign-side wealth is a documented
quasi-unit-root state whose one-step image ran to $+20.6\%$ against a $\pm8\%$ box.
\textbf{A box must be able to contain the image of its own points.}

\paragraph{The $P_\D$ band.} Symmetric at $\pm12\%$, because the period map has no
solution above roughly $P_\D+15\%$. An asymmetric $(0.20,0.12)$ pair was tried and is
\emph{worse} (six failures): it shifts the box centre off the steady state, so the
steady state stops being a collocation node --- and a solution that is not exact at
its own rest point is a poor starting point for everything else.

\subsection{Invariance and the rotated box}

Ideally the box would be \emph{one-step invariant}: the image of every point under
the transition law would lie inside it. Measured at the steady state, the Jacobian of
the deposit-obligation block is
\begin{equation}
J=\frac{\partial\bigl(P'_\D,P'_\F\bigr)}{\partial\bigl(P_\D,P_\F\bigr)}
=\begin{pmatrix}1.076 & -0.310\\ 2.016 & -1.252\end{pmatrix},
\qquad
\mathrm{eig}(J)=\{0.766,\,-0.942\} .
\label{eq:Pjacobian}
\end{equation}
The dynamics are stable --- both eigenvalues are inside the unit circle --- but $J$
is strongly non-normal (singular values $2.61$ and $0.28$). Invariance of an
axis-aligned box requires $|J|\,b\le b$ componentwise for some positive half-width
vector $b$, where $|J|$ is the entrywise absolute value. Here
$\rho(|J|)=1.96>1$, so by Perron--Frobenius \textbf{no axis-aligned box centred on
the steady state is one-step invariant, at any bandwidth}. On the eigenbasis, by
contrast, the map is diagonal with $|\lambda|<1$ and every box is invariant.

The code therefore supports collocating on rotated coordinates $z=R(\bS-\bar\bS)$
(Bocola's $V$-transform), with the box drawn on $z$ while every public method still
speaks natural coordinates:

\begin{lstlisting}[caption={Rotated box coordinates. \code{state\_grid.py}, \code{SmolyakGrid}.}]
def _fwd(self, x):
    # NATURAL -> ROTATED BOX COORDINATES (identity when rot is None).
    x = np.atleast_2d(x)
    return x if self.rot is None else (x - self.centre) @ self.rot.T

def clip(self, x):
    # PROJECT POINTS INTO THE BOX (SIMULATION USE; EXPLICIT, NOT SILENT).
    # Clipping happens in the ROTATED coordinates the box is drawn on, so a
    # rotated box projects onto its own faces, not onto an axis-aligned hull.
    return self._bwd(np.clip(self._fwd(x), self.lo, self.hi))
\end{lstlisting}

\warn{The rotation is correct in theory and \emph{measures worse}: $11$ of $41$
points fail and the worst residual is $1.2\times10^{-2}$, against zero failures and
$1.1\times10^{-14}$ for the axis-aligned box. The reason is instructive. A rotated
box is a parallelogram in natural coordinates, and its corners reach $P_\F$ states at
which the period map has no equilibrium at all. Invariance was simply not the binding
constraint; \emph{feasibility of the collocation points} was. The facility is kept,
tested, and off by default (\code{ROTATE\_P = False} in \code{main.py}).}

\subsection{The alternative: ergodic-set grids}

When the box approach fails outright --- as it does at $\ell=2$, where $40$--$50\%$ of
the collocated corners are infeasible --- the structural fix is to stop collocating on
a box at all. The \code{eds.py} module implements the Maliar--Maliar
$\varepsilon$-distinguishable-set method: simulate the model to obtain the ergodic
cloud, keep a well-spaced subset by farthest-point sampling, and fit by
complete-polynomial least squares on that cloud. The solver then only ever visits
states the model actually reaches. \code{EDSGrid} is a drop-in replacement exposing
the same \code{points / n / d / fit / eval / basis / clip} interface, so
\code{RuleSet}, \code{point\_map} and the time-iteration driver are reused unchanged
--- an illustration of why the grid abstraction is worth the indirection.

%%=====================================================================
\clearpage
\section*{Part 2. From theory to the residual function}
\addcontentsline{toc}{section}{Part 2. From theory to the residual function}
%%=====================================================================

\section{Assembling the residual}
\label{sec:residual}

\subsection{What gets parameterised}

Seventeen functions are approximated on the grid, in each of two default regimes ---
$34$ coefficient vectors in total. They fall into three groups:

\begin{table}[htbp]
\centering\small
\caption{The approximated rules. Source: \code{decision\_rules.py}.}
\label{tab:rules}
\begin{tabular}{@{}l>{\raggedright\arraybackslash}p{3.5cm}p{7.0cm}@{}}
\toprule
Group & Members & Role \\
\midrule
\code{SOLVE7}   & \code{N\_D}, \code{N\_F}, \code{Kp\_D}, \code{Kp\_F},
                  \code{rdep\_D}, \code{rdep\_F}, \code{p}
                & The pointwise Newton unknowns $\bx$: solved at every node from the
                  seven market-clearing residuals. \\
\addlinespace
\code{DERIVED} (valuations) & \code{alpha\_D}, \code{alpha\_F},
                  \code{Q\_bD}, \code{Q\_bF}, \code{r\_wc\_D}, \code{r\_wc\_F}
                & Banker valuations $\varphi_X$, bond prices $q^{X}$ and the lending
                  rate: \emph{read off} the Euler recursions given the frozen
                  continuation, not solved for. \\
\addlinespace
\code{DERIVED} (household) & \code{C\_D}, \code{C\_F}, \code{A\_D}, \code{A\_F}
                & Aggregate consumption and deposits, from the budget identity and
                  quantity clearing. \\
\bottomrule
\end{tabular}
\end{table}

The split between ``solved'' and ``read off'' is a design choice with real
consequences, discussed in \S\ref{sec:knot}: solving for all seventeen jointly at
each node would be a $17\times17$ non-linear system per point instead of $7\times7$,
and --- more importantly --- would reintroduce a simultaneity between the bond price
and next period's state that the frozen-continuation device removes.

\subsection{The seven equilibrium conditions}

At a node $(d,\bS)$, given a frozen continuation rule set $\Gamma$, the residual
function is
\begin{equation}
R:\R^{7}\to\R^{7},
\qquad
R\bigl(\bx;\,d,\bS,\Gamma\bigr)=0,
\label{eq:R-def}
\end{equation}
with components as in Table~\ref{tab:seven}. Each row states the economics, the
residual as coded, and --- crucially --- the \emph{normalisation}.

\begin{table}[htbp]
\centering\small
\caption{The residual vector. Source: \code{point\_map.py}, \code{point\_residuals}.}
\label{tab:seven}
\begin{tabular}{@{}clp{6.4cm}c@{}}
\toprule
\# & Condition & Residual as coded & Pins \\
\midrule
1 & Capital Euler, $\D$ &
$\bigl(\E[\Omega_\D(r^{k\prime}_\D-r_\D)]-\lambda_K\mu_\D\bigr)\big/\E[\Omega_\D]$ & $K'_\D$\\
2 & Capital Euler, $\F$ &
$\bigl(\E[\Omega_\F(r^{k\prime}_\F-r_\F)]-\lambda_K\mu_\F\bigr)\big/\E[\Omega_\F]$ & $K'_\F$\\
3 & Labour market, $\D$ &
$\bigl(\chi_\D N_\D^{1/\nu}-w_\D/P^{c}_\D\bigr)\big/\bigl(w_\D/P^{c}_\D\bigr)$ & $N_\D$\\
4 & Labour market, $\F$ &
$\bigl(\chi_\F N_\F^{1/\nu}-w_\F/P^{c}_\F\bigr)\big/\bigl(w_\F/P^{c}_\F\bigr)$ & $N_\F$\\
5 & Deposit Euler, $\D$ &
$x_\D^{-\sigma}\big/\bigl(\tilde\beta_\D\,\E[(1+r^{c\prime}_\D)x_\D'^{-\sigma}]\bigr)-1$ & $r_\D$\\
6 & Deposit-UIP &
$(1+r_\D)-(1+r_\F)\E[p']/p+\kappa_{\rm nfa}(P_\D-P_\F)/\bar Y_\D$ & $r_\F$\\
7 & Goods market, $\D$ &
$\bigl(Y_\D-P^{c}_\D C_\D-I_\D-NX_\D-G_\D\bigr)\big/\bar Y_\D$ & $p$\\
\bottomrule
\end{tabular}
\end{table}

\trick{\textbf{Normalisation is not cosmetic.} A Newton solver measures convergence
by $\|R\|_\infty$, so residuals in wildly different units make the tolerance
meaningless and the step direction dominated by whichever equation happens to be
measured in the largest units. Every residual above is therefore rendered
dimensionless: Euler conditions are divided by the discount kernel $\E[\Omega]$ so
they read in \emph{return} units, the labour conditions by the real wage, and the
goods market by steady-state output. Note the choice \emph{not} made: dividing the
capital Euler by $\lambda_K\bar\mu$ instead of $\E[\Omega]$ would amplify it roughly
$130\times$ relative to the others and wrecks the solver step.}

\subsection{Structure of one residual evaluation}

A single call to \code{point\_residuals} performs the following, in order. This is
the innermost loop of the entire solve, executed hundreds of times per node per
sweep, so its structure is worth reading carefully.

\begin{enumerate}
\item \textbf{Unpack} the trial vector $\bx$ and the state $\bS$.
\item \textbf{Firms and capital, current period.} Evaluate the production and
      capital blocks: output, the frictionless wage, $mpk$, Tobin's
      $q$, investment and capital-producer rents, in each country.
\item \textbf{Government.} Apply the fiscal rule to the surviving stock, compute the
      coupon, and forward-integrate the debt stock to $B'$.
\item \textbf{Balance sheets.} Compute asset payoffs $\mathcal{X}_X$, gross wealth
      $n^{g}_X=\mathcal{X}_X-P_X$, end-of-period assets $\mathcal{A}_X$, net worth
      $n_X$, deposits, and the next-period obligation $P'_X$.
\item \textbf{Continuations.} Build the next-period state vector
      $\bS'$ for every quadrature node, clip it into the box, and evaluate
      \emph{all} rules in each regime (\code{cont.eval\_all}).
\item \textbf{Branch kernels.} Form $\Lambda^{(d')}$, $\varphi^{(d')}$ and hence
      $\Omega^{(d')}$ on each branch; take the three-branch expectations.
\item \textbf{Banker block.} The closed-form multiplier $\mu_X$, the franchise value
      $\varphi_X$, bond prices $q^{X}$, the lending rate $r^{wc}_X$.
\item \textbf{Household block.} Dividends, income, consumption from the budget, the
      deposit Euler.
\item \textbf{Trade.} Import demand, net exports.
\item \textbf{Assemble} the seven residuals and return them together with a
      dictionary \code{out} of every derived object, which becomes the next iterate's
      rule values.
\end{enumerate}

The returned pair \code{(res, out)} is the reason one function serves both purposes:
the solver consumes \code{res}, and the sweep stores \code{out}.

\subsection{The occasionally-binding constraint inside the residual}

The multiplier is not a Newton unknown. It is computed \emph{inside} the residual
evaluation from Proposition~1 of the companion document,
\begin{equation}
\mu_{X}=\min\left\{\max\left\{1-\frac{\E[\Omega_X]\,(1+r_X)\,n_X}{\lambda\mathcal{A}_X},\;0\right\},\;
\bar\mu\right\},
\qquad \bar\mu=0.95,
\label{eq:mu-closed-form}
\end{equation}
and the capital Euler (rows 1--2 of Table~\ref{tab:seven}) is then the equation that
pins $K'$. This is the recursive-stable form of the same KKT system that the older
path-based solver imposed as a Fischer--Burmeister complementarity,
\begin{equation}
\Psi^{FB}(\mu,\mathrm{slack})=\mu+\mathrm{slack}-\sqrt{\mu^{2}+\mathrm{slack}^{2}}=0
\;\Longleftrightarrow\;
\mu\ge0,\ \mathrm{slack}\ge0,\ \mu\cdot\mathrm{slack}=0 .
\label{eq:FB}
\end{equation}
Both encode \eqref{eq:kkt} exactly. The difference is numerical: the FB form is a
smooth-except-at-the-origin residual added to a stacked system, whereas
\eqref{eq:mu-closed-form} is an explicit, bounded formula evaluated pointwise.

\begin{lstlisting}[caption={The KKT switch (hard) inside a smoothed cap. \code{point\_map.py}.}]
    # closed-form mu in [0, mu_cap]: the lower max is Bocola's KKT switch (hard, as in
    # his residual_model.m -- it IS the complementarity, not a numerical guard); the
    # upper cap keeps alpha = E_Om R/(1-mu) finite when a default-regime net worth n
    # goes negative, and is SMOOTHED because unlike the KKT switch it is a guard that
    # would otherwise plant a plateau in a fitted object.
    _MU_CAP = 0.95
    mu_D = float(_smin(max(1.0 - E_Om_D * (1.0 + rdep_D) * n_D / lev_D, 0.0),
                       _MU_CAP, _GUARD_EPS))
\end{lstlisting}

\warn{The distinction in that comment is the single most important line of
implementation philosophy in the solver. The \emph{lower} \code{max} is economics ---
it is the complementarity condition itself --- and must stay hard, because smoothing
it would blur the regime boundary the model is about. The \emph{upper} \code{min} is
a numerical guard against a diverging recursion, and must be smoothed, because a hard
cap plants a plateau inside an object that is subsequently fitted by a global
polynomial. Kinks that are economics stay; kinks that are numerics get smoothed. See
\S\ref{sec:smoothing}.}

\section{Computing expectations: quadrature over a continuum and a fork}
\label{sec:quadrature}

\subsection{Gauss--Hermite quadrature}

Every conditional expectation in the model is an integral against the standard normal
density,
\begin{equation}
\E_t\bigl[g(\varepsilon')\bigr]
=\int_{-\infty}^{\infty}g(\varepsilon)\,\frac{1}{\sqrt{2\pi}}e^{-\varepsilon^{2}/2}\,d\varepsilon .
\label{eq:expectation-integral}
\end{equation}
Gaussian quadrature approximates this by a weighted sum at optimally-chosen nodes,
\begin{equation}
\E_t\bigl[g(\varepsilon')\bigr]\;\approx\;\sum_{m=1}^{M}w_m\,g(\varepsilon_m),
\label{eq:gh-rule}
\end{equation}
where the nodes $\varepsilon_m$ are the roots of the degree-$M$ orthogonal polynomial
for the weight function and the weights $w_m$ are chosen to integrate low-order
polynomials exactly. The defining property is remarkable: an $M$-node Gaussian rule
is exact for all polynomials of degree $\le 2M-1$. With $M=5$ nodes the rule
integrates polynomials up to degree nine exactly, which for smooth integrands is
far more accurate than a $5$-point Monte Carlo draw (whose error falls only as
$M^{-1/2}$).

\subsection{The probabilists' versus physicists' trap}

There are two conventions for Hermite polynomials, and confusing them silently
rescales the shock.

\begin{itemize}
\item \textbf{Physicists' Hermite} $H_n$ is orthogonal with respect to
      $e^{-x^{2}}$, so its rule approximates
      \[\textstyle\int g(x)\,e^{-x^{2}}\,dx\;\approx\;\sum_m \tilde w_m\,g(x_m).\]
      Recovering an $N(0,\sigma^{2})$ expectation then requires the map
      $\varepsilon=\sqrt{2}\,\sigma x_m$ and weights $\tilde w_m/\sqrt{\pi}$.
\item \textbf{Probabilists' Hermite} $He_n$ is orthogonal with respect to
      $e^{-x^{2}/2}$. Its nodes are already in standard-deviation units, so
      $\varepsilon=\sigma x_m$ with weights normalised to sum to one.
\end{itemize}
This code uses the probabilists' rule, so the natural map $s'=\text{mean}+\sigma_s x_m$
is exact:

\begin{lstlisting}[caption={Quadrature nodes. \code{point\_map.py}, \code{gh\_nodes}.}]
def gh_nodes(n=7):
    # GAUSS-HERMITE NODES/WEIGHTS FOR ONE STANDARD-NORMAL INNOVATION.
    # hermegauss is the PROBABILISTS' rule (weight exp(-x^2/2)), so the nodes are
    # already in standard-deviation units and s' = mean + sigma*node is exact.
    # (Bocola's GaussHermite.m returns the PHYSICISTS' rule and then uses the same
    # sigma*node map, which silently rescales his innovation by 1/sqrt(2).)
    x, w = np.polynomial.hermite_e.hermegauss(n)
    return x, w / w.sum()
\end{lstlisting}

\warn{The parenthetical is a real, documented discrepancy with the reference
implementation, and it is worth understanding as a cautionary tale. Applying the
$\sigma\times\text{node}$ map to physicists' nodes implicitly divides the innovation
standard deviation by $\sqrt2$. The reference model's \emph{reported} parameter is
$\sigma_s=0.63$, but the model it actually solves contains
$\sigma_s^{\text{eff}}=0.63/\sqrt2=0.4455$. To reproduce those numbers exactly one
must set \code{cal["sigma\_s"] = 0.4455} here --- reproducing the solved model, not
the reported parameter. Two lines of quadrature convention change the effective
volatility of the driving process by $29\%$.}

\subsection{Mixing a continuous shock with a discrete fork}

The expectation is not over $\varepsilon'$ alone. Next period the sovereign either
does not default ($d'=0$), or defaults and the central-bank backstop is honoured, or
defaults and the backstop is reneged. The measure is therefore a product of the
Gauss--Hermite rule with a three-point discrete distribution, giving $3M$ evaluation
nodes:
\begin{equation}
\boxed{\;
\E_t\bigl[g\bigr]=\sum_{m=1}^{M}w_m
\Bigl\{\underbrace{(1-\pi^{d})}_{w^{\rm nd}}g\bigl(0,\bS'_m\bigr)
+\underbrace{\pi^{d}\varkappa}_{w^{\rm hon}}g^{\rm hon}\bigl(\bS'_m\bigr)
+\underbrace{\pi^{d}(1-\varkappa)}_{w^{\rm ren}}g\bigl(1,\bS'_m\bigr)\Bigr\}\;}
\label{eq:three-branch}
\end{equation}
with $\pi^{d}=\pi^{d}(s_t)$ the current logistic default probability and $\varkappa$
the priced backstop-activation probability. The implementation folds the branch
probabilities directly into the quadrature weight vectors, so a branch expectation is
a pair of dot products:

\begin{lstlisting}[caption={Branch weights and the expectation operators. \code{point\_map.py}.}]
    w_nd, w_def = wq * (1.0 - pd), wq * pd
    a_tpi = float(cal.get("tpi_activation", 0.0))
    w_hon, w_ren = w_def * a_tpi, w_def * (1.0 - a_tpi)

    def _E(v0, v1):
        # TWO-BRANCH EXPECTATION (no-default vs default) -- F-safe legs.
        return float(np.dot(w_nd, v0) + np.dot(w_def, v1))

    def _E3(v0, vh, vr):
        # THREE-BRANCH EXPECTATION: no-default / default-honoured / default-reneged.
        return float(np.dot(w_nd, v0) + np.dot(w_hon, vh) + np.dot(w_ren, vr))
\end{lstlisting}

\trick{\textbf{Exact nesting by construction.} Setting $\varkappa=0$ makes
$w^{\rm hon}\equiv0$ in exact floating-point arithmetic, so
\eqref{eq:three-branch} collapses to the two-branch case with no residual
contamination --- not approximately, but bit-for-bit. Likewise $\pi^{d}\equiv0$ makes
$w^{\rm def}\equiv0$ and recovers the risk-neutral model exactly. This is what makes
the nesting regression test (\code{tests/test\_recursive\_nesting.py}) a meaningful
check rather than a tolerance comparison, and it is why the branch probabilities are
multiplied into the weights rather than branched on with an \code{if}.}

\subsection{Evaluating the continuation}

Each of the $M$ quadrature nodes needs a full next-period state. These are assembled
as a single $(M\times7)$ array and the rules are evaluated on all of them in one
vectorised call --- one basis-matrix construction and one matrix multiply per regime,
rather than $M$ scalar interpolations:

\begin{lstlisting}[caption={Vectorised continuation evaluation. \code{point\_map.py}, \code{\_cont}.}]
    def _cont(d_next):
        # CONTINUATION RULE VALUES + NEXT-PERIOD CAPITAL RETURN, REGIME d_next.
        Sn = np.empty((m, 7))
        Sn[:, IK_D] = Kp_D; Sn[:, IK_F] = Kp_F
        Sn[:, IP_D] = Pp_D; Sn[:, IP_F] = Pp_F
        Sn[:, IB_D] = Bp_D
        Sn[:, IS] = s_next
        Sn[:, IZ] = Z_next
        r = cont.eval_all(d_next, cont.grid.clip(Sn))
        # guard continuation outputs before they enter fractional powers / sqrt
        # (a deep default-regime iterate can push N, p, C negative -> NaN)
        for k in ("N_D", "N_F"):
            r[k] = np.maximum(r[k], 0.05)
        for k in ("C_D", "C_F"):
            r[k] = np.maximum(r[k], 1e-3)
        r["p"] = np.maximum(r["p"], 1e-2)
\end{lstlisting}

Note the ordering: only $s'$ varies across quadrature nodes; the endogenous states
$K'$, $P'$, $B'$ are common to all $M$ of them, because they are chosen at $t$ and
are therefore deterministic given $\bx$. The array is broadcast accordingly.

\subsection{The honoured branch as a blend}

The backstop-honoured continuation is not a separately solved regime --- there is no
third coefficient set. It is a convex blend of the two solved regimes' rules,
governed by the ``real shield'' parameter $\varsigma$:
\begin{equation}
g^{\rm hon}(\bS')=(1-\varsigma)\,g\bigl(1,\bS'\bigr)+\varsigma\,g\bigl(0,\bS'\bigr),
\label{eq:blend}
\end{equation}
implemented as a dictionary comprehension over all rules at once:

\begin{lstlisting}[caption={The honoured branch. \code{point\_map.py}.}]
    # honoured continuation = d'=1 recession blended toward d'=0 by `shield`
    rh = {k: (1.0 - shield) * r1[k] + shield * r0[k] for k in r1}
\end{lstlisting}

This is an approximation --- a genuinely separate backstop regime would require a
third solve --- but a disciplined one: at $\varsigma=1$ the honoured branch is exactly
the no-default economy, at $\varsigma=0$ it is exactly the default economy, and
intermediate values interpolate. The baseline sets $\varsigma=0.5$ because at
$\varsigma=1$ full activation neutralises risk entirely and tips the barely-binding
constraint into a slack-branch slip.

\subsection{Quadrature order}

The solve uses $M=5$ nodes (\code{N\_GH = 5} in \code{recursive\_experiment.py}). The
order is \emph{stamped onto the solved rule set}:

\begin{lstlisting}[caption={Recording the quadrature order. \code{decision\_rules.py}, \code{RuleSet}.}]
        # quadrature order the rules were SOLVED under; time_iteration stamps it and
        # every reader (IRFs, decompositions, accuracy) uses it, so a solve at n_gh=5
        # is never read back at n_gh=7 and charged the difference as approximation error
        self.n_gh = None
\end{lstlisting}

\trick{This is a small piece of bookkeeping that prevents a subtle and embarrassing
class of bug. If rules are solved with $M=5$ and then read back with $M=7$, the
residuals evaluated at the solution will not be zero --- and the difference is pure
quadrature mismatch, not approximation error. Every downstream reader (impulse
responses, decompositions, the accuracy report) takes its order from
\code{rules.n\_gh}, so the reported accuracy is the accuracy of the object that was
actually solved.}

%%=====================================================================
\clearpage
\section*{Part 3. Code architecture and implementation}
\addcontentsline{toc}{section}{Part 3. Code architecture}
%%=====================================================================

\section{Architecture}
\label{sec:architecture}

\subsection{Module map}

The solver is deliberately layered, with the economics kept solver-agnostic. The
call structure of a production run is:

\begin{center}\small
\begin{tabular}{@{}llp{7.4cm}@{}}
\toprule
Layer & Module & Responsibility \\
\midrule
Driver        & \code{main.py}                     & Steady state $\to$ TFP $\to$ risk $\to$ TPI experiments \\
              & \code{recursive\_experiment.py}    & Grid choice, warm-start chain, homotopy, IRF readers \\
              & \code{tpi\_recursive\_experiment.py} & Backstop-activation comparison \\
\addlinespace
Solver        & \code{recursive\_main.py}          & \code{time\_iteration}, \code{\_sweep}, \code{solve\_point} \\
              & \code{point\_map.py}               & \code{point\_residuals}: the residual $R(\bx;d,\bS,\Gamma)$ \\
\addlinespace
Approximation & \code{state\_grid.py}              & \code{SmolyakGrid}, box construction, the $s$ process \\
              & \code{decision\_rules.py}          & \code{RuleSet}: coefficients, transforms, evaluation \\
              & \code{eds.py}                      & Ergodic-set grid (drop-in alternative) \\
\addlinespace
Economics     & \code{blocks/*.py}                 & Bank, firms, capital, trade, government, household \\
              & \code{config/*.py}                 & Calibration and the two-stage steady-state solve \\
\addlinespace
Diagnostics   & \code{accuracy.py}                 & Euler errors on the ergodic set, box escapes \\
              & \code{recursive\_residual.py}      & Per-point $\mu$/slack/residual maps \\
              & \code{output\_decomposition.py}    & Factor decomposition of $d\log Y$ \\
\bottomrule
\end{tabular}
\end{center}

The dependency direction is strict: \code{blocks/} knows nothing about grids,
quadrature or solvers; \code{point\_map.py} calls \code{blocks/} but does not know
about time iteration; \code{recursive\_main.py} knows about grids and residuals but
contains no economics. This is what allowed the entire perfect-foresight solver stack
to be deleted without touching a single economic block, and what allows
\code{EDSGrid} to substitute for \code{SmolyakGrid} unchanged.

\subsection{The \code{RuleSet} abstraction}

A \code{RuleSet} holds, for each of the $17$ rules and each of the two regimes, both
the point values on the grid and the fitted coefficients. Its interface is small:

\begin{lstlisting}[caption={The rule container. \code{decision\_rules.py}.}]
class RuleSet:
    # COEFFICIENTS AND POINT VALUES FOR EVERY RULE IN BOTH REGIMES.
    def __init__(self, grid):
        self.grid = grid
        self.vals = {k: [np.empty(grid.n), np.empty(grid.n)] for k in ALL_RULES}
        self.coef = {k: [None, None] for k in ALL_RULES}
        self.n_gh = None

    def set_values(self, name, d, values, weights=None, ridge=0.0):
        # SET POINT VALUES FOR ONE RULE IN ONE REGIME AND REFIT ITS COEFFICIENTS.
        self.vals[name][d] = np.asarray(values, dtype=float).copy()
        y = to_fit(name, self.vals[name][d])
        if weights is None and ridge == 0.0:
            self.coef[name][d] = self.grid.fit(y)
        else:
            self.coef[name][d] = self.grid.fit_weighted(y, weights, ridge)

    def eval_all(self, d, x):
        # EVALUATE EVERY RULE IN REGIME d AT POINTS x -> DICT OF ARRAYS.
        B = self.grid.basis(x)
        return {k: from_fit(k, B @ self.coef[k][d]) for k in ALL_RULES}
\end{lstlisting}

\trick{\code{eval\_all} builds the basis matrix \emph{once} and reuses it across all
seventeen rules. Since the basis evaluation is $O(n\cdot d\cdot k_{\max})$ and the
subsequent projection is a single matrix-vector product per rule, this makes
evaluating all rules barely more expensive than evaluating one. Given that
\code{\_cont} is called two or three times per residual evaluation, and the residual
is called dozens of times per Newton solve at each of $518$ node-regime pairs per
sweep, this single design choice is worth an order of magnitude.}

\subsection{Log-space collocation}

Rules are \emph{stored} in levels but \emph{fitted} in transformed space:

\begin{lstlisting}[caption={The fit transform. \code{decision\_rules.py}.}]
LOG_RULES   = frozenset({"N_D", "N_F", "Kp_D", "Kp_F", "p", "alpha_D", "alpha_F",
                         "Q_bD", "Q_bF", "C_D", "C_F", "A_D", "A_F"})
GROSS_RULES = frozenset({"rdep_D", "rdep_F", "r_wc_D", "r_wc_F"})

def to_fit(name, v):
    # LEVELS -> THE QUANTITY ACTUALLY FITTED BY THE CHEBYSHEV COLLOCATION.
    if name in GROSS_RULES:
        return np.log(np.maximum(1.0 + np.asarray(v, dtype=float), _FIT_FLOOR))
    if name in LOG_RULES:
        return np.log(np.maximum(np.asarray(v, dtype=float), _FIT_FLOOR))
    return np.asarray(v, dtype=float)

def from_fit(name, y):
    # FITTED QUANTITY -> LEVELS (the inverse of to_fit).
    if name in GROSS_RULES:
        return np.exp(np.clip(y, -50.0, 50.0)) - 1.0
    if name in LOG_RULES:
        return np.exp(np.clip(y, -50.0, 50.0))
    return y
\end{lstlisting}

Three things are going on here, and each is a distinct trick.

\begin{enumerate}
\item \textbf{Positivity by construction.} Hours, capital, consumption, prices and
      valuations must be positive. A polynomial fitted to levels can go negative
      between nodes --- especially near a kink, where Gibbs oscillation overshoots ---
      and a negative $N$ entering $N^{1/\nu}$ produces a NaN that propagates through
      every subsequent expectation. Fitting $\log x$ and exponentiating makes
      negativity \emph{impossible}, which removed a whole family of hard clips that
      previously existed only to protect positivity.
\item \textbf{No plateaus to represent.} Those hard clips were themselves harmful:
      a clip is a plateau, a plateau is a kink, and the interpolant then has to
      represent it. Removing them improves the fit for a second, independent reason.
\item \textbf{Rates handled separately.} A net interest rate may legitimately be
      negative, so $\log r$ is inadmissible. Rates are therefore carried as the
      \emph{gross} rate $1+r$, which is positive --- the same device Bocola uses when
      he parameterises policies as $\exp(\bar x+\gamma)$.
\end{enumerate}

The $\pm50$ clamp on the exponent is a last-resort finiteness guard: $\exp$ overflows
to infinity around $709$, and one infinity poisons every expectation it enters. The
range $e^{\pm50}$ spans $10^{-22}$ to $5\times10^{21}$ and is unreachable by any
admissible economic value, so the clamp never binds on a converged solution --- it
only keeps a diverging transient iterate finite long enough to be rejected.

\section{Mathematics-to-code map}
\label{sec:mapping}

Table~\ref{tab:map} indexes every mathematical object in Parts 1--2 to the exact file,
function and variable that implements it.

\begin{longtable}{@{}p{4.6cm}p{4.4cm}p{5.4cm}@{}}
\caption{Mathematics to code.}\label{tab:map}\\
\toprule
Mathematical object & Code location & Names \\
\midrule
\endfirsthead
\multicolumn{3}{@{}l}{\emph{Table~\ref{tab:map}, continued}}\\
\toprule
Mathematical object & Code location & Names \\
\midrule
\endhead
\bottomrule
\endfoot
$T_k(x)$, recurrence \eqref{eq:cheb-recurrence}
  & \code{state\_grid.py} & \code{chebyshev\_basis\_1d(x, max\_deg)} \\
Chebyshev extrema \eqref{eq:cheb-extrema}
  & \code{state\_grid.py} & \code{\_level\_points(i)} \\
New points $\Delta_i$
  & \code{state\_grid.py} & \code{\_new\_points(i)} \\
New degrees $\mathcal{K}_i$
  & \code{state\_grid.py} & \code{\_new\_degrees(i)} \\
Index set $\mathcal{I}(\ell)$ \eqref{eq:multi-index}
  & \code{state\_grid.py} & \code{\_multi\_indices(d, mu, mu\_vec)} \\
Sparse grid $\mathcal{H}(\ell)$ \eqref{eq:smolyak-grid}
  & \code{state\_grid.py} & \code{SmolyakGrid.points}, \code{.points\_unit} \\
Basis set $\mathcal{B}(\ell)$ \eqref{eq:smolyak-basis}
  & \code{state\_grid.py} & \code{SmolyakGrid.degrees} \\
Collocation matrix $\Phi$ \eqref{eq:phi-1d}
  & \code{state\_grid.py} & \code{\_Phi}, \code{\_basis\_unit(u)} \\
$c=\Phi^{-1}y$ \eqref{eq:fit-1d}
  & \code{state\_grid.py} & \code{fit(values)} via \code{lu\_solve} \\
Weighted ridge fit
  & \code{state\_grid.py} & \code{fit\_weighted(values, w, ridge)} \\
Box map $u(\bS)$ \eqref{eq:box-map}
  & \code{state\_grid.py} & \code{to\_unit}, \code{from\_unit} \\
Box construction \eqref{eq:bands}
  & \code{state\_grid.py} & \code{build\_state\_box(...)} \\
Rotation $z=R(\bS-\bar\bS)$
  & \code{state\_grid.py} & \code{\_fwd}, \code{\_bwd}, \code{rot}, \code{centre} \\
Clip \eqref{eq:clip}
  & \code{state\_grid.py} & \code{clip(x)}, \code{outside(x)} \\
Eigenbasis probe \eqref{eq:Pjacobian}
  & \code{recursive\_main.py} & \code{p\_block\_rotation(...)} \\
\addlinespace
State vector $\bS$ \eqref{eq:state}
  & \code{state\_grid.py} & \code{STATE\_NAMES}; indices \code{IK\_D...IZ} \\
Unknowns $\bx$ \eqref{eq:policy-vector}
  & \code{decision\_rules.py} & \code{SOLVE7} \\
Read-off rules
  & \code{decision\_rules.py} & \code{DERIVED} \\
$\hat g(\bS;c)$ \eqref{eq:parameterisation}
  & \code{decision\_rules.py} & \code{RuleSet.eval}, \code{.eval\_all} \\
Log transform
  & \code{decision\_rules.py} & \code{to\_fit}, \code{from\_fit} \\
\addlinespace
Residual $R(\bx;d,\bS,\Gamma)$ \eqref{eq:R-def}
  & \code{point\_map.py} & \code{point\_residuals(S, d, x, cont, ...)} \\
Closed-form $\mu$ \eqref{eq:mu-closed-form}
  & \code{point\_map.py} & \code{mu\_D}, \code{mu\_F}, \code{lev\_D}, \code{lev\_F} \\
Franchise value $\varphi$
  & \code{point\_map.py} & \code{alpha\_D\_new}, \code{alpha\_F\_new} \\
Bond prices $q^{X}$
  & \code{point\_map.py} & \code{Q\_bD\_new}, \code{Q\_bF\_new} \\
Lending spread $\lambda\mu/\E[\Omega]$
  & \code{point\_map.py} & \code{r\_wc\_D}, \code{wedge\_sp\_D} \\
GH nodes \eqref{eq:gh-rule}
  & \code{point\_map.py} & \code{gh\_nodes(n)} \\
Branch weights \eqref{eq:three-branch}
  & \code{point\_map.py} & \code{w\_nd}, \code{w\_hon}, \code{w\_ren} \\
Branch expectations
  & \code{point\_map.py} & \code{\_E(v0,v1)}, \code{\_E3(v0,vh,vr)} \\
Continuation states
  & \code{point\_map.py} & \code{\_cont(d\_next)} \\
Honoured blend \eqref{eq:blend}
  & \code{point\_map.py} & \code{rh}, \code{shield} \\
Smoothed guards
  & \code{point\_map.py} & \code{\_smax}, \code{\_smin}, \code{\_sclip} \\
\addlinespace
Point solve
  & \code{recursive\_main.py} & \code{solve\_point(S, d, cont, ..., x0)} \\
One sweep
  & \code{recursive\_main.py} & \code{\_sweep(rules, cont, ...)} \\
Time iteration
  & \code{recursive\_main.py} & \code{time\_iteration(rules, ...)} \\
Warm-start chain, homotopy
  & \code{recursive\_experiment.py} & \code{solve\_recursive(...)} \\
IRF reader
  & \code{recursive\_experiment.py} & \code{read\_at(...)} \\
Euler errors
  & \code{accuracy.py} & \code{simulate}, \code{euler\_errors}, \code{accuracy\_report} \\
EGM / distribution kernels
  & \code{fast\_kernels.py} & \code{hh\_backward}, \code{dist\_forward} \\
\end{longtable}

%%=====================================================================
\clearpage
\section*{Part 4. Computational tricks and optimisations}
\addcontentsline{toc}{section}{Part 4. Computational tricks}
%%=====================================================================

\section{The outer algorithm: time iteration}
\label{sec:timeiter}

\subsection{Why time iteration rather than one big Newton}

Given the parameterisation, the collocation conditions \eqref{eq:collocation} are a
system of
\begin{equation}
\underbrace{7}_{\text{equations}}\times\underbrace{259}_{\text{nodes}}
\times\underbrace{2}_{\text{regimes}}=3{,}626
\label{eq:system-size}
\end{equation}
non-linear equations in as many unknowns. One could hand that to a Newton solver
directly. The code does not, and the reason is the occasionally-binding constraint:
a monolithic Newton must construct a $3626\times3626$ Jacobian by finite differences,
and \emph{any} node whose perturbed evaluation crosses the KKT switch
$\{\mu=0\}$ contaminates an entire column. With thousands of nodes, some node is
almost always at the kink. The earlier \code{make\_collocation\_residual} approach
did exactly this and was superseded.

Instead the solver uses \textbf{time iteration} (a Coleman operator): treat the
previous iterate's rules as a \emph{frozen continuation} $\Gamma$, and for each node
independently solve the small $7\times7$ system $R(\bx;d,\bS,\Gamma)=0$. Formally,
define the operator
\begin{equation}
\mathcal{T}\bigl[\Gamma\bigr](d,\bS)\;=\;\bigl\{\bx^{\star},\;\text{derived}\bigr\}
\quad\text{where}\quad R\bigl(\bx^{\star};d,\bS,\Gamma\bigr)=0 ,
\label{eq:coleman}
\end{equation}
and iterate $\Gamma_{k+1}=\mathcal{T}[\Gamma_k]$ to a fixed point. Three advantages
follow immediately:

\begin{enumerate}
\item each node's solve is $7\times7$, so a kink at one node is confined to that node;
\item nodes are independent, so a failure is \emph{localisable} and can be handled
      (\S\ref{sec:masking}) rather than poisoning the global step;
\item the fixed point of $\mathcal{T}$ is exactly the collocation solution, so
      nothing is given up.
\end{enumerate}

\subsection{The sweep}

One sweep loops over regimes and nodes, solving each and collecting the results:

\begin{lstlisting}[caption={One time-iteration sweep. \code{recursive\_main.py}, \code{\_sweep}.}]
def _sweep(rules, cont, cal, ss, sproc, regimes, no_default, n_gh, keep_tol=1e-3):
    # ONE TIME-ITERATION SWEEP: SOLVE EVERY POINT, RETURN NEW RULE VALUE ARRAYS.
    # A point whose solve does not clear (fn > keep_tol) RETAINS the previous
    # iterate's values -- a failed corner must never poison the continuation.
    n = rules.grid.n
    new = {k: {d: np.empty(n) for d in regimes} for k in STORE()}
    wt = {d: np.ones(n) for d in regimes}     # per-point fit weight (0 = failed corner)
    worst, n_fail = 0.0, 0
    for d in regimes:
        for i in range(n):
            S = rules.grid.points[i]
            x0 = np.array([rules.vals[k][d][i] for k in SOLVE7])
            x, out, fn = solve_point(S, d, cont, cal, ss, sproc, x0,
                                     no_default=no_default, n_gh=n_gh, x_ss=x_ss)
            worst = max(worst, fn if np.isfinite(fn) else 1e3)
            if (not np.isfinite(fn)) or fn > keep_tol:   # retain old values, mask fit
                n_fail += 1
                wt[d][i] = 0.0
                for k in STORE():
                    new[k][d][i] = rules.vals[k][d][i]
            else:
                for j, k in enumerate(SOLVE7):
                    new[k][d][i] = x[j]
                for k in DERIVED:
                    new[k][d][i] = out[k]
    return new, worst, n_fail, wt
\end{lstlisting}

\subsection{Damping and the convergence test}

The updated values are not adopted outright. They are damped,
\begin{equation}
\Gamma_{k+1}=\varpi\,\mathcal{T}[\Gamma_k]+(1-\varpi)\,\Gamma_k,
\qquad \varpi=0.25,
\label{eq:damping}
\end{equation}
and only then refitted. Damping is what keeps the fixed-point iteration stable when
the operator is not a contraction --- which is exactly the situation near the
constraint boundary, where a small change in the continuation can move $\mu$ across
zero and produce a large change in policies. The cost is speed: with $\varpi=0.25$ the
per-sweep contraction is at best $1-\varpi=0.75$, so a cold start needs $70$--$90$
sweeps.

The exit test is deliberately a \emph{conjunction}:

\begin{lstlisting}[caption={Damping, refit and the two-part exit test. \code{recursive\_main.py}, \code{time\_iteration}.}]
        for k in STORE():
            for d in regimes:
                old = rules.vals[k][d]
                upd = damp * new[k][d] + (1.0 - damp) * old
                change = max(change, np.max(np.abs(upd - old))
                             / (np.max(np.abs(old)) + 1e-8))
                weights = (np.where(wt[d] > 0.5, 1.0, fit_fw) if fit_mask else None)
                rules.set_values(k, d, upd, weights=weights, ridge=fit_ridge)
        if change < tol and worst < 1e-6:
            return True, it + 1, worst, n_fail
    return False, max_it, worst, n_fail
\end{lstlisting}

\warn{The condition is \code{change < tol AND worst < 1e-6}, and \code{n\_fail} is
part of the return contract. The reason is a failure mode that occurred and went
unnoticed: a sweep can look perfectly converged on \code{change} while a quarter of
the grid is \emph{frozen} on its previous values and never clearing. The rules stop
moving --- because the failing nodes stop moving by construction --- and a
change-only test declares success on a solution that does not satisfy its own
equations at $25\%$ of its collocation points. The caller must be able to see this,
so the failure count is returned, and a non-converged stage prints a warning
\emph{even when \code{verbose} is off}, because the TPI and decomposition experiments
call the solver with \code{verbose=False} and silence there is exactly how a run
reporting failure into \code{\_} went unnoticed.}

\section{The inner solver}
\label{sec:newton}

\subsection{Powell's hybrid method}

Each node's $7\times7$ system is solved by \code{scipy.optimize.root(method="hybr")},
a modified Powell hybrid method (MINPACK's \code{hybrd}). It is a trust-region
Newton variant: it takes a Gauss--Newton step when the model is trusted and a
gradient-descent step otherwise, which makes it far more robust than pure Newton to
the poor initial guesses that occur early in the iteration.

The Jacobian is not supplied analytically --- deriving $\partial R/\partial\bx$
through a three-branch quadrature over interpolated continuations would be an
enormous and fragile undertaking --- so MINPACK builds it by forward differences with
step $\sqrt{\epsilon_{\text{mach}}}\approx1.49\times10^{-8}$ relative to each
component. \textbf{This finite-difference step size is the single number that governs
every smoothing decision in \S\ref{sec:smoothing}}: a guard whose kink is narrower
than $1.5\times10^{-8}$ is invisible to the solver in a bad way --- the difference
quotient straddles it and returns a meaningless slope.

\subsection{Dual starts and graceful failure}

\begin{lstlisting}[caption={The point solve. \code{recursive\_main.py}, \code{solve\_point}.}]
def solve_point(S, d, cont, cal, ss, sproc, x0, no_default=False, n_gh=7, x_ss=None):
    # SOLVE THE 7 MARKET-CLEARING UNKNOWNS AT ONE POINT (FROZEN CONTINUATION).
    # hybr from the warm start (the common case: one cheap solve near the fixed
    # point); a single fallback from the SS guess only if that misses.
    def f(x):
        try:
            return point_residuals(S, d, x, cont, cal, ss, sproc,
                                   n_gh=n_gh, no_default=no_default)[0]
        except (ValueError, RuntimeError, FloatingPointError):
            return np.full(7, 10.0)

    sol = root(f, x0, method="hybr", tol=1e-12)
    best = (sol.x, np.max(np.abs(sol.fun)))
    if best[1] > 1e-9 and x_ss is not None:
        sol2 = root(f, x_ss, method="hybr", tol=1e-12)
        if np.max(np.abs(sol2.fun)) < best[1]:
            best = (sol2.x, np.max(np.abs(sol2.fun)))
\end{lstlisting}

Three techniques appear in these few lines.

\paragraph{Exceptions as penalties.} An economically infeasible trial vector --- one
that makes the Jermann inversion demand a negative bracket, say --- raises rather than
returning a NaN. The wrapper catches it and returns a large constant residual
$(10,\dots,10)$, which steers the trust region away from that region instead of
crashing the sweep. The alternative, returning NaN, would poison the Jacobian.

\paragraph{Warm starting.} The initial guess \code{x0} is the \emph{previous
iterate's solved values at this same node}. Near the fixed point this is an excellent
guess and the solve converges in a handful of iterations. This is the single largest
speed factor in the whole algorithm.

\paragraph{A second start, only if needed.} If the warm start misses ($\|R\|>10^{-9}$),
one retry from the steady-state guess is attempted and the better of the two is kept.
Attempting it unconditionally would double the cost for no benefit; not attempting it
at all leaves nodes stranded when the warm start is in the wrong basin.

Finally, if \emph{both} evaluations raise, the function returns a sentinel
$\|R\|=10^{3}$ rather than propagating the exception, so \code{\_sweep} can mask the
node and continue.

\section{Continuation and homotopy}
\label{sec:homotopy}

\subsection{Breaking the simultaneity knot}
\label{sec:knot}

There is a genuine circularity in the model: the current bond price $q^{\D}_t$ enters
the government's financing need, which determines $B'$, which is part of next
period's state, which determines the continuation values that price $q^{\D}_t$. The
same loop exists for the franchise value $\varphi$ and for the working-capital rate
$r^{wc}$ (whose loan quantity depends on the wage, which depends on $r^{wc}$, which
depends on $\mu$, which depends on the balance sheet).

The solver cuts all three knots the same way: \textbf{take the circular object from
the frozen previous iterate.}

\begin{lstlisting}[caption={Knot-breaking with frozen valuations. \code{point\_map.py}.}]
    # current-period bank valuations = FROZEN previous-iterate rules at this state
    # (breaks the S'<->Q_b knot; coincide at the fixed point)
    Sm = np.atleast_2d(S)
    Q_bD_cur    = float(cont.eval("Q_bD",   d, Sm)[0])
    alpha_D_cur = float(cont.eval("alpha_D", d, Sm)[0])
    r_wc_D_cur  = float(cont.eval("r_wc_D", d, Sm)[0])
\end{lstlisting}

\trick{This is the defining idea of time iteration applied \emph{within} the period,
and its correctness argument is simple: at the fixed point, the frozen value and the
freshly computed value coincide by construction, so the converged solution satisfies
the true simultaneous system. Away from the fixed point they differ, but that
difference is just another dimension of the iteration error, which the convergence
test measures. What it buys is that the $7\times7$ system stays $7\times7$ and stays
solvable.}

Note the internal consistency requirement this creates: the labour residual uses the
\emph{contemporaneous} $r^{wc}$ (computed fresh from $\mu$), while the loan
\emph{quantity} uses the frozen one. The two coincide at the fixed point; using the
frozen value in both places would leave the labour condition unenforced.

\subsection{Warm-start chaining across regimes}

The default regime $d=1$ cannot be cold-started: it is a deep recession far from the
steady state, and a cold Newton from steady-state values does not converge at any
node. The driver therefore chains:

\begin{enumerate}
\item Solve $d=0$ at $\pi^{d}\equiv0$ (the risk-neutral economy) from a steady-state
      cold start. This is the easiest problem in the family.
\item Copy the converged $d=0$ values into the $d=1$ slots as an initial guess.
\item Solve $d=1$ by homotopy in the recovery rate (below).
\item Solve both regimes \emph{jointly} with risk priced, which is the first stage in
      which the two regimes see each other.
\end{enumerate}

\subsection{Homotopy in the recovery rate}

Step 3 is a continuation method. Rather than jumping to the calibrated haircut, the
solver walks the recovery rate down in four steps, re-solving at each:

\begin{lstlisting}[caption={Haircut homotopy. \code{recursive\_experiment.py}, \code{solve\_recursive}.}]
    # d=1 (post-default) via a HAIRCUT HOMOTOPY: warm-start from d=0, then lower
    # the recovery from ~1 (=no default, identical to d=0) down to the calibrated
    # value in steps, re-solving each -- so the deep recession is reached by
    # continuation, not cold-started (which does not converge).
    for k in STORE_RULES:
        rules.set_values(k, 1, rules.vals[k][0].copy())
    for rec in (0.85, 0.70, 0.55, rec_target):
        cal["recovery_rate_D"] = rec
        ok1, it1, w1, f1 = time_iteration(rules, cal, ss, sproc, regimes=(1,),
                                          no_default=False, damp=0.25, tol=1e-6,
                                          max_it=80, n_gh=N_GH)
\end{lstlisting}

Formally, let $\varrho\in[\varrho_{\text{target}},1]$ index a family of problems
$\mathcal{P}(\varrho)$, with $\mathcal{P}(1)$ the trivial no-haircut problem whose
solution is known (it is the $d=0$ solution) and $\mathcal{P}(0.45)$ the one wanted.
The homotopy traces the solution path
$\varrho\mapsto\Gamma^{\star}(\varrho)$ by using the solution at each step as the
initial guess for the next. It works whenever the path is continuous and the steps
are small enough that each solution lies in the next problem's basin of attraction.

\trick{Homotopy is the standard remedy when a non-linear solve fails from a cold
start but the problem embeds in a family with a known easy member. The recovery rate
is the natural parameter here precisely because $\varrho=1$ makes the default branch
\emph{identical} to the no-default branch --- a perfect known solution at the end of
the path, not merely an approximate one.}

\section{Robustness of the fit}
\label{sec:masking}

\subsection{Failed nodes must not enter the fit}

Some collocation nodes have no equilibrium at all --- in the default regime, corners
that combine a deep haircut with already-low net worth are genuinely infeasible.
Since the fit \eqref{eq:fit-1d} is a \emph{global} least-squares/interpolation
problem, a single garbage value at one node moves the polynomial \emph{everywhere},
including at the ergodic centre where the answer is read. Measured: one saturated
node moves the interpolant at the ergodic centre by about $26\%$ and can drive it
negative.

Two defences are layered. First, a failed node retains its previous values, so the
value entering the fit is at least economically sensible. Second, the fit itself can
be switched from exact square collocation to a masked, ridge-regularised weighted
least squares:

\begin{lstlisting}[caption={Robust weighted ridge fit. \code{state\_grid.py}, \code{fit\_weighted}.}]
def fit_weighted(self, values, w=None, ridge=0.0):
    # RIDGE-REGULARISED WEIGHTED LEAST-SQUARES FIT: the fix for global-fit corner
    # poisoning at mu=2. w in [0,1] per point (0 = point EXCLUDED from the fit, so
    # an unsolvable/frozen corner cannot leak into the coefficients); the ridge
    # penalises high-degree coefficients (damps the Gibbs wiggle at the kink) with
    # a small uniform floor for invertibility when points are masked.
    values = np.asarray(values, dtype=float)
    if w is None and ridge == 0.0:
        return lu_solve(self._lu, values)
    w = np.ones(self.n) if w is None else np.asarray(w, dtype=float)
    A = self._Phi.T * w                                    # Phi^T @ diag(w)
    lhs = A @ self._Phi                                    # Phi^T W Phi
    diag_scale = float(np.mean(np.diag(lhs))) + 1e-12
    td = self.degrees.sum(axis=1).astype(float)            # total degree per basis fn
    reg = diag_scale * (ridge * td / max(td.max(), 1.0) + 1e-6)
    lhs[np.diag_indices_from(lhs)] += reg
    return np.linalg.solve(lhs, A @ values)
\end{lstlisting}

The estimator is
\begin{equation}
c=\Bigl(\Phi^{\top}W\Phi+\Lambda\Bigr)^{-1}\Phi^{\top}W\,y,
\qquad
\Lambda_{bb}=\bar\Phi\left(\rho\,\frac{|k_b|_1}{\max_b|k_b|_1}+10^{-6}\right),
\label{eq:ridge}
\end{equation}
with $|k_b|_1=\sum_j k_{b,j}$ the total degree of basis function $b$. Two design
details are worth noting.

\paragraph{Degree-graded ridge.} The penalty is proportional to the total degree, so
low-order (economically meaningful) terms are barely penalised while high-order terms
--- exactly those that produce Gibbs oscillation near a kink --- are damped hard. A
uniform ridge would shrink the level and slope of the rule, which is not the problem.

\paragraph{Soft rather than hard masking.} The default failure weight is $0.1$, not
$0$. Hard masking (weight zero) collapses the fit when a sweep has many failures ---
$\Phi^{\top}W\Phi$ loses rank --- so failed corners are kept as weak anchors: they
still constrain the fit enough to keep it full-rank and sane, but their poisoning is
cut roughly tenfold. The $10^{-6}$ uniform floor guarantees invertibility regardless.

\section{Smoothing, guards and bounds}
\label{sec:smoothing}

\subsection{The smoothing primitives}

Wherever a bound must be imposed on a quantity that is subsequently \emph{fitted} or
that enters a finite-difference Jacobian, the hard operator is replaced by its smooth
counterpart:
\begin{align}
\mathrm{smax}(x,\underline{x};\epsilon)
&=\underline{x}+\tfrac12\Bigl[(x-\underline{x})
   +\sqrt{(x-\underline{x})^{2}+\epsilon^{2}}\Bigr],
\label{eq:smax}\\[2pt]
\mathrm{smin}(x,\bar x;\epsilon)
&=\bar x-\tfrac12\Bigl[(\bar x-x)
   +\sqrt{(\bar x-x)^{2}+\epsilon^{2}}\Bigr].
\label{eq:smin}
\end{align}

\begin{lstlisting}[caption={Smoothing primitives. \code{point\_map.py}.}]
def _smax(x, floor, eps=1e-3):
    # SMOOTH MAX(x, floor) -- differentiable everywhere so the FD-Jacobian stays valid.
    return floor + 0.5 * ((x - floor) + np.sqrt((x - floor) ** 2 + eps ** 2))

def _sclip(x, lo, hi, eps=None):
    # SMOOTH clip: the guards below must not be plateaus. A hard np.clip on a value
    # that is then FITTED puts a kink inside the box, and one saturated node moves the
    # interpolant at the ergodic centre by ~26% and can drive it negative (measured).
    return _smin(_smax(x, lo, eps or _GUARD_EPS), hi, eps or _GUARD_EPS)
\end{lstlisting}

\subsection{\texorpdfstring{Calibrating $\epsilon$: the two-sided argument}
                          {Calibrating epsilon: the two-sided argument}}

The smoothing scale is not arbitrary. It is squeezed between two requirements.

\paragraph{Upper bound (bias).} Away from the bound, $\mathrm{smax}$ differs from
$\max$ by approximately $\epsilon^{2}/(4\,\text{gap})$, where ``gap'' is the distance
to the bound. Since the steady state must remain an exact rest point of the recursive
system to within the acceptance tolerance, $\epsilon$ must be small enough that this
bias is negligible there. At $\epsilon=10^{-5}$ the induced steady-state error is
about $10^{-11}$ --- three orders of magnitude below acceptance.

\paragraph{Lower bound (visibility).} The smoothing exists to hide the kink from the
finite-difference Jacobian, whose step is $\approx1.5\times10^{-8}$
(\S\ref{sec:newton}). If $\epsilon$ were comparable to or smaller than that step, the
difference quotient would straddle the smoothed region and see the kink anyway. At
$\epsilon=10^{-5}$ the smoothing is roughly three decades wider than the FD step.

\begin{equation}
\underbrace{1.5\times10^{-8}}_{\text{FD step}}\;\ll\;
\underbrace{\epsilon=10^{-5}}_{\code{\_GUARD\_EPS}}\;\ll\;
\underbrace{\sqrt{4\,g\,\tau}}_{\text{bias limit}} .
\label{eq:eps-window}
\end{equation}

The net-worth floor keeps its own larger $\epsilon=10^{-3}$, because it operates on a
quantity of order $n_{ss}$ rather than on an $O(1)$ ratio.

\subsection{Inventory of guards}

\begin{table}[htbp]
\centering\small
\caption{Every bound in the residual, and its status.}
\label{tab:guards}
\begin{tabular}{@{}>{\raggedright\arraybackslash}p{2.9cm}l>{\raggedright\arraybackslash}p{7.5cm}@{}}
\toprule
Guard & Value & Status and rationale \\
\midrule
KKT switch $\max\{\cdot,0\}$ on $\mu$ & --- &
\textbf{Hard.} It is the complementarity condition, i.e.\ economics. Smoothing it
would blur the regime boundary. \\
$\mu$ cap & $0.95$ &
\textbf{Smoothed.} Keeps $\varphi=\E[\Omega]R/(1-\mu)$ finite when default-regime net
worth goes negative. \\
$\varphi$ cap & $40$ &
\textbf{Smoothed.} The $\varphi$ recursion has slope $\beta(1-f)R/(1-\mu)$, which
exceeds one once $\mu>1-\beta(1-f)R\approx0.04$; above that the cap is what arrests a
divergent fixed point, not cosmetics. Tunable via \code{cal["alpha\_cap"]}. \\
Net-worth floor & $0.15\,n_{ss}$ &
\textbf{Smoothed}, $\epsilon=10^{-3}$. Bocola's own $\max\{\cdot,0.65\}$ device.
Inactive in the ergodic region ($\code{nw\_floor\_frac}=0$ reproduces the baseline);
catches only deep default corners. \\
Bond price bounds & $[0.2,\,1.2]$ &
\textbf{Smoothed.} Numerical backstops, never binding in the ergodic region (the
reference solution bottoms at $0.639$). \\
Consumption bounds & $[0.3,\,2.0]\times C_{ss}$ &
\textbf{Smoothed}, with $\epsilon$ scaled by $C_{ss}$ so the relative tightness
matches the $O(1)$ guards. \\
Continuation floors & $N\ge0.05$, $C\ge10^{-3}$, $p\ge10^{-2}$ &
\textbf{Hard} --- but applied to continuation \emph{inputs} before fractional powers,
not to fitted outputs, so no plateau is collocated. \\
Jermann bracket & $>0$ &
\textbf{Raises.} An infeasible investment rate throws \code{ValueError}, which the
solver wrapper converts into a penalty residual. \\
Exponent clamp & $\pm50$ &
\textbf{Hard.} Finiteness only; unreachable by admissible values. \\
\bottomrule
\end{tabular}
\end{table}

\warn{The cap interaction in row three is a real bug that was found and fixed, and it
illustrates why guards must be designed jointly. An older hard clip of $\varphi$ at
$8.0$ bound \emph{first} --- at $\mu\approx0.76$ --- and bit exactly where the
net-worth floor puts the deep default corner: $n=0.15\,n_{ss}$ implies
$\mu\approx0.85$ and $\varphi\approx13$, and the reference solution legitimately
reaches $\varphi/\bar\varphi=13.6$. Two guards intended to be slack were in fact
contradicting each other and truncating a region the model genuinely visits.}

\section{Vectorisation and low-level performance}
\label{sec:performance}

\subsection{Where the time goes}

A production risk-stage solve is roughly $90$ minutes: $259$ nodes $\times$ $2$
regimes $\times$ $\sim$$100$ sweeps $\times$ $\sim$$30$ residual evaluations per
Newton solve $\approx1.5$ million residual evaluations, each containing two or three
continuation evaluations of $17$ rules at $5$ quadrature nodes. The optimisations
below target that innermost product.

\subsection{Vectorisation}

\begin{itemize}
\item \textbf{Basis matrices, not scalar interpolations.} \code{eval\_all} builds one
      $(M\times n)$ basis matrix and applies it to all $17$ coefficient vectors
      (\S\ref{sec:architecture}).
\item \textbf{One LU factorisation for all fits.} $\Phi$ is factorised at grid
      construction; every one of the $34$ fits per sweep is a back-substitution.
\item \textbf{Quadrature as dot products.} Branch probabilities are folded into the
      weight vectors so an expectation is \code{np.dot}, with no Python-level loop
      over nodes or branches.
\item \textbf{Broadcast continuation states.} Only $s'$ varies across quadrature
      nodes; the endogenous states are assigned by broadcast into the $(M\times7)$
      array.
\item \textbf{Single \code{bincount} in the distribution step.} The Young lottery
      scatter is done with one \code{np.bincount} over flattened $(a,e)$ indices
      rather than a per-state \code{np.add.at} --- roughly $4\times$ faster.
\end{itemize}

\subsection{JIT compilation with an exact fallback}

The two hot loops of the heterogeneous-agent block --- EGM backward induction and the
distribution forward iteration --- are inherently sequential and cannot be vectorised
away. They are compiled with numba:

\begin{lstlisting}[caption={JIT with a transparent fallback. \code{fast\_kernels.py}.}]
try:
    from numba import njit
    HAVE_NUMBA = True
except ImportError:                                   # pragma: no cover
    HAVE_NUMBA = False
    def njit(*args, **kwargs):
        # NO-OP STAND-IN FOR numba.njit WHEN NUMBA IS UNAVAILABLE.
        def wrap(f):
            return f
        return wrap

@njit(cache=True)
def hh_backward(a_grid, Pi_T, r_path, y_path, c_terminal, beta, sigma, a_min, vN_path):
    # EGM BACKWARD INDUCTION OVER THE PATH (MIRRORS household.egm_step EXACTLY).
\end{lstlisting}

\trick{Three details make this safe. (i) The decorator is \emph{stubbed} when numba
is absent, so the module imports and runs identically without the dependency ---
slowly, but correctly. (ii) The kernels replicate the numpy reference operation for
operation, and the equivalence is regression-tested
(\code{tests/test\_fast\_kernels.py}), so the fast path can never silently diverge
from the reference semantics. (iii) \code{cache=True} persists the compiled
artefacts, which matters because spawned worker processes would otherwise each pay
the JIT compilation cost.}

Note in particular that \code{np.searchsorted} is hand-written as an explicit binary
search inside the JIT kernel, because numba's coverage of numpy's search routines
does not extend to the exact call signature used; the hand-rolled version is what
makes the kernel compile in \code{nopython} mode.

\subsection{The spawn guard}

On macOS, Python's \code{multiprocessing} uses the \emph{spawn} start method, which
re-imports the parent module in every worker. Any unguarded top-level code therefore
runs once per worker. This produced a real and initially baffling symptom --- ``SS
solved'' printed nine times --- and the rule is now absolute: every standalone script
that may fan out is wrapped in

\begin{lstlisting}
if __name__ == "__main__":
    main()
\end{lstlisting}

The recursive solver's pointwise solves are serial, so the hazard is currently
latent; the guard is kept regardless, since the modules are importable and the cost of
the guard is zero.

\section{Diagnosing the solution}
\label{sec:accuracy}

\subsection{Why residuals at the nodes are not enough}

By construction, collocation drives $R$ to zero \emph{at the nodes}. Reporting that
$\max_b\|R(\bS_b)\|=10^{-14}$ therefore says almost nothing about solution quality ---
it is close to a tautology. The meaningful question is how large the residual is
where the model actually spends its time, which is generally \emph{not} at the
collocation nodes.

\subsection{Euler errors on the ergodic set}

The diagnostic simulates the solved rules forward under drawn innovations and, at
every visited state, evaluates the seven equilibrium conditions \emph{at the
rule-implied policy} --- with no re-solving. The distribution of
$\log_{10}|R|$ per equation is the accuracy report.

\begin{lstlisting}[caption={Simulating the solved rules. \code{accuracy.py}, \code{simulate}.}]
    for t in range(T + burn):
        Sm = np.atleast_2d(S)
        x = np.array([float(rules.eval(k, 0, Sm)[0]) for k in SOLVE7])
        _, o = point_residuals(S, 0, x, rules, cal, ss, sproc, n_gh=ngh, ...)
        ...
        Sn = S.copy()
        Sn[0], Sn[1] = x[2], x[3]                       # K' = Kp
        Sn[2], Sn[3] = o["Pp_D"], o["Pp_F"]
        Sn[4] = o["Bp_D"]
        Sn[5] = ((1.0 - sproc["rho_s"]) * sproc["s_star"] + sproc["rho_s"] * S[5]
                 + sproc["sigma_s"] * rng.standard_normal())
        S = rules.grid.clip(Sn)[0]
\end{lstlisting}

Note that the endogenous states are advanced by the \emph{period map's own}
end-of-period values, so the simulated path is the model's, not the grid's.

\subsection{Box escapes}

The same routine reports how often, and by how much, the simulated path leaves the
collocation box --- per state dimension, as a fraction of box width:

\begin{lstlisting}[caption={Reporting box violations. \code{state\_grid.py} and \code{accuracy.py}.}]
def outside(self, x):
    # PER-DIMENSION BOX VIOLATION IN ROTATED COORDS, AS A FRACTION OF BOX WIDTH.
    z = self._fwd(x)
    return np.maximum(np.maximum(self.lo - z, z - self.hi), 0.0) / (self.hi - self.lo)
\end{lstlisting}

\trick{This turns the box-invariance discussion of \S\ref{sec:box} from theory into a
measurement. States are deliberately \emph{not} clipped before being recorded ---
leaving the box is a result to be reported, not a condition to be silently enforced.
If the ergodic path spends a material fraction of periods outside the box, the
reported impulse responses are partly extrapolation, and the bands must widen.}

\subsection{Reading impulse responses off the rules}

There is one subtlety in how results are read, and it is worth stating because the
obvious alternative is wrong.

\begin{lstlisting}[caption={Reading the solution, not re-solving it. \code{recursive\_experiment.py}, \code{read\_at}.}]
def read_at(rules, cal, ss, sproc, S):
    # READ THE CONVERGED RULES AT STATE S (binding branch), returning the implied
    # allocation + the point residual there (accuracy). Evaluating the period map
    # at the rules' OWN policy values stays on the binding branch -- re-solving
    # with a root finder can slip onto the nearby slack equilibrium at the barely-
    # binding SS. This is the standard way to read a global solution's IRF.
\end{lstlisting}

\warn{Re-solving at each impulse-response date is tempting --- it would drive the
labour-FOC residual to zero --- but it is wrong at a barely-binding steady state.
Measured at $\ell=1$, the Newton step slips onto the \emph{slack} equilibrium in $49$
of $60$ periods: $\mu_\D$ walks from $0.103$ to $0$ and the entire path turns
expansionary ($Y_\D$ $+2.5\%$ on impact against $-0.004\%$ for the rule read). The
residual seen when reading the rules is genuine approximation error, and the cure is
a finer grid --- not re-solving the current period, which changes which equilibrium
branch is selected.}

\section{\texorpdfstring{The known failure: why $\ell=2$ does not converge}
                       {The known failure: why level 2 does not converge}}
\label{sec:mu2}

Documenting what does \emph{not} work is part of the specification. The production
solve runs at $\ell=1$ with anisotropic refinement in $s$; the isotropic $\ell=2$
grid does not converge, and the reasons are understood.

\paragraph{Reason 1: the multiplier loses its bite.} At $\ell=2$ the sign of the
output response can flip. The mechanism is that a coarser or differently-weighted
representation lets $\E[\Omega]$ rise where it should fall, which by
\eqref{eq:mu-closed-form} drives $\mu\to0$ --- the constraint goes slack precisely
when it should bind. A separate study established that binding harder at the steady
state (raising $f$ and the target spread so that $\bar\mu=0.057$) fixes the sign, but
that this is a calibration change, not a numerical fix, and it was not adopted.

\paragraph{Reason 2: corner poisoning.} Even with the sign repaired, the joint
residual remains of order one. At $\ell=2$ some $40$--$50\%$ of collocated corners are
infeasible: a wide box is forced by the near-unit-root capital states, and the
post-default $d=1$ regime simply has no equilibrium at its cross-corners. Those
points cannot be fitted, tightened, or masked away in sufficient number --- with that
many failures, hard masking collapses the fit and soft masking still leaks.

\paragraph{The structural fix.} The corners must not be \emph{collocated} at all,
which is what the ergodic-set grid of \code{eds.py} implements: fit on the simulated
cloud, where every point is feasible by construction, using a low complete-degree
basis that cannot produce Gibbs oscillation. This is the documented next step rather
than a completed result.

\trick{The general lesson transfers to any projection solve of a model with an
occasionally-binding constraint and a discrete regime: \textbf{a global polynomial
basis is only as good as the feasibility of its collocation set}. Refining the grid
increases accuracy only if the added points are points at which the model has a
solution. Where they are not, refinement makes things strictly worse --- which is the
opposite of the intuition carried over from one-dimensional numerical analysis.}

%%=====================================================================
\clearpage
\appendix
\section{Reproducing the results}
%%=====================================================================

All commands are run from \code{code/global/} with plain \code{python3} (numpy,
scipy, matplotlib; numba optional).

\begin{center}\footnotesize
\begin{tabular}{@{}>{\raggedright\arraybackslash}p{6.9cm}>{\raggedright\arraybackslash}p{7.8cm}@{}}
\toprule
Command & What it does \\
\midrule
\code{python3 main.py} &
Full pipeline: steady state, TFP experiment, sovereign-risk pass-through, accuracy
report, TPI activation overlay. $\sim$2\,h. \\
\code{python3 -m}\newline
\code{\ \ solver\_recursive.}\newline
\code{\ \ recursive\_experiment} &
Sovereign-risk stage only. \\
\code{python3 -m}\newline
\code{\ \ solver\_recursive.}\newline
\code{\ \ tpi\_recursive\_experiment} &
OMT/TPI activation comparison at $0/50/100\%$. \\
\code{python3 tests/test\_state\_grid.py} &
Smolyak counts, collocation exactness, quadratic exactness, box coverage of the $s$
process, rotated-box round trip. \\
\code{python3 tests/test\_recursive\_nesting.py} &
Steady state is a rest point; $\pi^{d}=0$ nests the risk-neutral model. \\
\code{python3 tests/test\_fast\_kernels.py} &
numba/numpy kernel equivalence. \\
\bottomrule
\end{tabular}
\end{center}

The principal solver switches live at the top of \code{main.py}: \code{MU} (TFP grid
level), \code{RISK\_MU\_VEC} (the anisotropic risk grid), \code{NW\_FLOOR},
\code{ROTATE\_P}, \code{RUN\_TPI} and \code{TPI\_ACTIVATIONS}.

\section{Glossary of numerical terms}

\begin{center}\small
\begin{tabular}{@{}lp{9.6cm}@{}}
\toprule
Term & Meaning in this document \\
\midrule
Projection method & Replace an unknown function by a finite basis expansion and
impose the functional equation in a weighted-residual sense. \\
Collocation & The projection variant that sets the residual to zero exactly at $n$
chosen nodes. \\
Chebyshev extrema & The Gauss--Lobatto nodes \eqref{eq:cheb-extrema}; nested, endpoint-inclusive. \\
Smolyak grid & Sparse union of small tensor products under a total-level budget
\eqref{eq:multi-index}. \\
Approximation level $\ell$ & The Smolyak budget (\code{mu} in code); higher means more
points and higher degree. \\
Anisotropic level $\ell_j$ & Per-dimension budget (\code{mu\_vec}); concentrates
resolution on chosen states. \\
Time iteration & Fixed-point iteration on decision rules using the previous iterate as
a frozen continuation (a Coleman operator). \\
Frozen continuation & The previous iterate's rules, held fixed while the current
period is solved; coincides with the solution at the fixed point. \\
Homotopy / continuation & Solving a hard problem by tracing a path from an easy member
of a parameterised family. \\
Damping & Convex combination of the new and old iterates,
$\varpi\mathcal{T}[\Gamma]+(1-\varpi)\Gamma$. \\
Gauss--Hermite quadrature & Optimal node/weight rule for Gaussian integrals; exact for
polynomials of degree $\le2M-1$. \\
Ridge / Tikhonov & Penalised least squares; here graded by basis total degree to damp
high-order oscillation. \\
Gibbs oscillation & Spurious ringing of a global polynomial near a kink or
discontinuity. \\
Euler error & Residual of an equilibrium condition evaluated at the approximate policy;
the standard accuracy metric. \\
EDS grid & $\varepsilon$-distinguishable set: a well-spaced subsample of the simulated
ergodic cloud, used as collocation points. \\
\bottomrule
\end{tabular}
\end{center}

\end{document}
